% methods.tex — Online Methods for PWM Flagship Nature Paper
% This file is included via % methods.tex — Online Methods for PWM Flagship Nature Paper
% This file is included via % methods.tex — Online Methods for PWM Flagship Nature Paper
% This file is included via % methods.tex — Online Methods for PWM Flagship Nature Paper
% This file is included via \input{methods} from the main manuscript.

\section*{Online Methods}

%% =========================================================================
\subsection*{OperatorGraph Specification}
\label{sec:methods:operatorgraph}

\paragraph{Formal definition.}
The \opg{} intermediate representation encodes the forward physics of any
computational imaging modality as a directed acyclic graph (DAG)
$\mathcal{G} = (\mathcal{V}, \mathcal{E})$.
Each node $v_i \in \mathcal{V}$ wraps a \emph{primitive operator} and
implements two entry points: $\texttt{forward}(x) \to y$ and
$\texttt{adjoint}(y) \to x$, the latter defined only when the primitive
is linear.  Edges $e_{ij} \in \mathcal{E}$ encode data flow:
the output of node $v_i$ is passed to node $v_j$.
Each node additionally exposes a set of learnable parameters
$\boldsymbol{\theta}_i$ that may be perturbed during mismatch simulation
or optimized during calibration, as well as read-only metadata flags
(\texttt{is\_linear}, \texttt{is\_stochastic}, \texttt{is\_differentiable}).
The graph is stored as a declarative YAML specification
(\texttt{OperatorGraphSpec}) and compiled to an executable
\texttt{GraphOperator} object by the \texttt{GraphCompiler}.

\paragraph{Node types.}
Primitive operators fall into two categories:

\begin{itemize}
  \item \textbf{Linear operators.}  Convolution (\texttt{conv2d}), mask
    modulation (\texttt{mask\_modulate}), sub-pixel shift (\texttt{subpixel\_shift\_2d}),
    Radon transform (\texttt{radon\_fanbeam}),
    Fourier encoding (\texttt{fourier\_encode}), spectral dispersion
    (\texttt{spectral\_disperse}), Fresnel propagation
    (\texttt{fresnel\_propagate}), random projection
    (\texttt{random\_project}), and structured illumination
    (\texttt{sim\_modulate}).
    Each implements both \texttt{forward()} and \texttt{adjoint()}.
  \item \textbf{Nonlinear operators.}  Squared magnitude
    (\texttt{magnitude\_sq}), Poisson--Gaussian noise
    (\texttt{poisson\_gaussian}), saturation clipping
    (\texttt{saturation\_clip}), phase retrieval nonlinearity
    (\texttt{phase\_abs}), and detector quantization
    (\texttt{quantize}).
    These set \texttt{is\_linear = False} and raise
    \texttt{NotImplementedError} on \texttt{adjoint()}, except where a
    well-defined pseudo-adjoint exists (\eg{}, the identity adjoint for
    magnitude-squared in Gerchberg--Saxton-type algorithms).
\end{itemize}

\paragraph{Adjoint consistency and Lipschitz bounds as mathematical safety rails.}
Two structural constraints distinguish the \opg{} IR from standard deep-learning
forward models and prevent either nonlinear primitive family from becoming a
universal approximator.

First, \emph{adjoint consistency}: correctness of every linear primitive
is verified by a randomized dot-product test.  For a primitive $A$ with
forward map $A:\mathbb{R}^n \to \mathbb{R}^m$, we draw $x \sim \mathcal{N}(0,I_n)$
and $y \sim \mathcal{N}(0,I_m)$ and compute
\begin{equation}
  \delta = \frac{|\langle A^* y,\, x \rangle - \langle y,\, Ax \rangle|}
               {\max(|\langle A^* y,\, x \rangle|,\; \epsilon)}
  \label{eq:adjoint_check}
\end{equation}
where $\epsilon = 10^{-12}$ guards against division by zero.  The test
is repeated $n_{\text{trials}} = 5$ times with independent random draws;
the primitive passes if $\delta_{\max} < 10^{-6}$.  At the graph level,
a compiled \texttt{GraphOperator} composed entirely of linear nodes
executes the same test over the composed forward--adjoint chain.  A
\texttt{GraphAdjointCheckReport} records $n_{\text{trials}}$,
$\delta_{\max}$, and $\bar\delta$ for audit.  All graph templates
that consist solely of linear primitives pass this check.
Adjoint consistency is not merely a numerical correctness guarantee: it
is a \emph{physical faithfulness certificate}.  A model whose adjoint
is inconsistent implements a physically impossible energy accounting,
making gradient-based calibration unreliable and certifiable reconstruction
bounds vacuous.  Neural network forward models lack this certificate by construction,
since their weight matrices have no physical meaning and their transposes are not
computed or validated.

Second, \emph{Lipschitz bounds}: the two nonlinear primitive families,
Detect~$D$ and Transform~$\Lambda$, are each restricted to five canonical
function families with at most two scalar parameters.  This restriction
enforces bounded Lipschitz constants ($\mathrm{Lip}(\Lambda_k) \leq L$ for
the explicit $L$ derived in~\cite{yang2026fpt}) and prevents either primitive
from becoming a universal approximator in the sense of Cybenko~\cite{cybenko1989approximation}.
Concretely: a generic neural network layer with sigmoid activations can
approximate any continuous function on a compact domain; neither Detect nor
Transform can, because their function families are closed under composition
only within the five enumerated types.  This non-universality is the key
property that makes the 11-primitive library \emph{falsifiable}: if a
modality's forward physics cannot be represented within the restricted
families, the extension protocol (Main Text, ``Extension protocol'') fires
and a new primitive must be justified physically.  The combination of
adjoint consistency and bounded Lipschitz nonlinearity provides the
mathematical safety rails that make \opg{}-based reconstruction both
interpretable and certifiable---properties that standard physics-informed
neural networks~\cite{raissi2019pinn} achieve only approximately and
without modality-agnostic guarantees.

\paragraph{Graph compilation.}
The compiler executes a four-stage pipeline:
\begin{enumerate}
  \item \textbf{Validate.}  Confirm acyclicity via topological sort
    (Kahn's algorithm), verify that every \texttt{primitive\_id}
    exists in the global \texttt{PRIMITIVE\_REGISTRY}, reject
    duplicate \texttt{node\_id} values, and optionally verify shape
    compatibility along edges when a \texttt{canonical\_chain}
    metadata flag is set.
  \item \textbf{Bind.}  Instantiate each primitive with its parameter
    dictionary $\boldsymbol{\theta}_i$.
  \item \textbf{Plan forward.}  The topological sort yields a sequential
    execution plan $(v_{\pi(1)}, \ldots, v_{\pi(|\mathcal{V}|)})$.
  \item \textbf{Plan adjoint.}  For graphs where
    \texttt{all\_linear = True}, the adjoint plan reverses the
    topological order and applies each node's individual adjoint in
    sequence, implementing the chain rule $A^* = A_1^*
    \circ \cdots \circ A_{|\mathcal{V}|}^*$ for a composition $A = A_{|\mathcal{V}|}
    \circ \cdots \circ A_1$.  For graphs containing nonlinear nodes,
    the adjoint plan is not generated, and any call to
    \texttt{adjoint()} raises \texttt{NotImplementedError} at runtime.
\end{enumerate}
The compiled \texttt{GraphOperator} is serializable to JSON and hashable
via SHA-256 for provenance tracking in RunBundle manifests.

\paragraph{Template library.}
The \texttt{graph\_templates.yaml} registry contains 170 templates
spanning all five carrier families. Of these, 12 modalities carry full
end-to-end Scenario~I--IV correction validation across all five carrier
families (Supplementary Table~S3). Representative modalities by carrier family include:

\begin{itemize}
  \item \textbf{Optical and X-ray photons:}  CASSI, SPC, CACTI, structured
    illumination microscopy (SIM), confocal, light-sheet, holography,
    ptychography, Fourier ptychographic microscopy (FPM), optical coherence
    tomography (OCT), lensless imaging, light field, integral imaging,
    neural radiance fields (NeRF), Gaussian splatting, fluorescence
    lifetime imaging (FLIM), diffuse optical tomography (DOT),
    phase retrieval, X-ray computed tomography (CT), and cone-beam CT
    (CBCT).
  \item \textbf{Electrons:}  Electron diffraction, electron backscatter
    diffraction (EBSD), electron energy loss spectroscopy (EELS), and
    electron holography.
  \item \textbf{Nuclear spins (MRI):}  Functional MRI (fMRI), diffusion-weighted
    MRI (DW-MRI), and magnetic resonance spectroscopy (MRS).
  \item \textbf{Acoustic waves:}  Ultrasound B-mode, Doppler ultrasound,
    shear-wave elastography, sonar, and photoacoustic tomography
    (combines optical excitation with acoustic detection).
  \item \textbf{Particle-beam probes:}  Neutron tomography, proton CT,
    and muon tomography.  Each decomposes into existing primitives with
    no extension required: neutron CT follows the X-ray CT DAG
    ($P \!\to\! \Lambda \!\to\! D$, Beer--Lambert attenuation with nuclear
    cross-section $\mu_n$); proton CT inserts a Bethe--Bloch
    stopping-power stage as Transform~$\Lambda$ (\#11); muon tomography
    uses Scatter~$R$ (\#10) for multiple Coulomb scattering and
    POCA vertex reconstruction (Supplementary Note~S24).
\end{itemize}

\paragraph{Fidelity levels.}
Each template is parameterized by a fidelity level that controls the
degree of physical realism in the simulated forward model:

\begin{description}
  \item[Level~1 (Linear, shift-invariant):]  The forward model is a
    linear, spatially uniform operator---the simplest approximation,
    suitable for initial diagnostics and rapid prototyping.
  \item[Level~2 (Linear, shift-variant):]  Spatially varying operator
    parameters (e.g.\ non-uniform illumination, position-dependent PSF,
    multi-coil sensitivity maps in MRI).  Adds a modality-appropriate
    noise model (Poisson shot noise plus Gaussian read noise for
    photon-counting modalities, Rician noise for MRI, Poisson for CT).
  \item[Level~3 (Nonlinear, ray/wave-based):]  Includes nonlinear
    effects such as beam hardening, phase wrapping, and scattering,
    captured by the Transform~$\Lambda$ primitive.
    Perturbation families and ranges are specified in
    \texttt{mismatch\_db.yaml}.
  \item[Level~4 (Full-wave / Monte Carlo):]  Complete physical
    simulation including wave-optical propagation, spatially varying
    aberrations, detector nonlinearities, and environmental drift.
    Currently implemented for holography and ptychography.
\end{description}

\noindent All 11 primitives in the library $\Blib$ apply uniformly across every fidelity level; the levels organize simulation complexity, not theorem scope.


%% =========================================================================
\subsection*{Triad Decomposition Formalization}
\label{sec:methods:triad}

The \triad{} asserts that the quality of any computational imaging
reconstruction is bounded by three fundamental gates.  Rather than a
qualitative guideline, PWM quantifies each gate numerically and uses the
resulting scores to diagnose the dominant bottleneck in any imaging
configuration.

\paragraph{Gate~1 (Recoverability).}
Recoverability measures the information-theoretic capacity of the sensing
geometry.  We quantify it via the \emph{effective compression ratio}
$r = m / n$, where $m$ is the number of independent measurements and $n$
the dimension of the signal.  The \texttt{compression\_db.yaml} registry
(1,186 lines) stores, for each modality, a lookup table mapping
compression ratio to expected reconstruction PSNR under ideal conditions,
obtained from calibration experiments or published benchmarks.  Each entry
carries a \texttt{provenance} field citing the source (paper DOI, internal
experiment ID, or theoretical formula).  Additional recoverability
indicators include the effective rank of the measurement matrix (estimated
via randomized SVD for large operators), the dimension of the null space,
and the restricted isometry property (RIP) constant where analytically
tractable (\eg{}, for Gaussian random projections in SPC).

\paragraph{Gate~2 (Carrier Budget).}
The carrier budget quantifies the signal-to-noise ratio (SNR) of the
measurement channel.  The photon-budget module consumes the
\texttt{photon\_db.yaml} registry which stores, per modality,
a deterministic photon model parameterized by source power, quantum
efficiency, exposure time, and detector characteristics.  It
classifies the noise regime into one of three categories:
\emph{shot-limited} (Poisson-dominated, SNR $\propto \sqrt{N_{\text{photon}}}$),
\emph{read-limited} (Gaussian read noise dominates, SNR $\propto
N_{\text{photon}} / \sigma_{\text{read}}$), and \emph{dark-current-limited}
(long exposures where dark current accumulation dominates).  The output is
a \texttt{PhotonReport} containing the estimated SNR in decibels, the noise
regime classification, per-element photon count, and a feasibility verdict
(\texttt{sufficient}, \texttt{marginal}, or \texttt{insufficient}).

\paragraph{Gate~3 (Operator Mismatch).}
Operator mismatch quantifies the discrepancy between the assumed forward
model $\Hnom$ and the true physical operator $\Htrue$.  The
mismatch scoring module consults \texttt{mismatch\_db.yaml}
which catalogs, for each modality, the set of mismatch parameters
(spatial shifts, rotational offsets, dispersion errors, PSF deviations,
coil sensitivity errors, center-of-rotation offsets, \etc{}), their
typical ranges, and available correction methods.  The mismatch severity
score $s \in [0, 1]$ is computed as the normalized $\ell_2$ distance
$\| \boldsymbol{\theta}_{\text{true}} - \boldsymbol{\theta}_{\text{nom}} \|
/ \| \boldsymbol{\theta}_{\text{range}} \|$, where
$\boldsymbol{\theta}_{\text{range}}$ is the per-parameter dynamic range
from the registry.  Sensitivity analysis $\partial \psnr / \partial
\theta_k$ is estimated via finite differences on the forward model.  The
output is a \texttt{MismatchReport} containing the severity score, the
dominant mismatch parameter, the recommended correction method, and the
expected PSNR gain from correction.

\paragraph{Gate binding determination.}
Given reconstruction results under the four-scenario protocol
(the Evaluation Protocol section below), PWM identifies the dominant gate by
comparing three cost terms:
\begin{align}
  C_{\text{mismatch}} &= \psnr_{\text{I}} - \psnr_{\text{II}} \label{eq:cost_mismatch} \\
  C_{\text{noise}}    &= \psnr_{\text{ideal}} - \psnr_{\text{noisy}} \label{eq:cost_noise} \\
  C_{\text{recover}}  &= \psnr_{\text{limit}} - \psnr_{\text{I}} \label{eq:cost_recover}
\end{align}
where $\psnr_{\text{I}}$ is the reconstruction PSNR under Scenario~I
(ideal operator), $\psnr_{\text{II}}$ under Scenario~II (mismatched
operator), $\psnr_{\text{noisy}}$ under the corresponding noisy condition,
and $\psnr_{\text{limit}}$ is the theoretical upper bound from the
compression table.  The dominant gate is $\arg\max_g C_g$.

\paragraph{TriadReport.}
For every diagnosis, \pwm{} produces a \triadreport{}: a Pydantic-validated
structured artifact comprising
\texttt{dominant\_gate} (enum: \texttt{recoverability},
\texttt{carrier\_budget}, \texttt{operator\_mismatch}),
\texttt{evidence\_scores} (three floats, one per gate),
\texttt{confidence\_interval} (float, 95\% CI width from bootstrap),
\texttt{recommended\_action} (string, \eg{} ``increase compression ratio''
or ``apply mismatch correction''), and
\texttt{parameter\_sensitivities} (dictionary mapping each mismatch
parameter name to its $\partial \psnr / \partial \theta_k$ value).
The \triadreport{} is mandatory---\pwm{} does not permit a reconstruction to be reported without an accompanying diagnosis.

\paragraph{Recovery ratio.}
We define the \emph{recovery ratio}
\begin{equation}
  \recoveryratio = \frac{\psnr_{\text{IV}} - \psnr_{\text{II}}}
                        {\psnr_{\text{I}} - \psnr_{\text{II}}}
  \label{eq:recovery_ratio}
\end{equation}
which lies in $[0, 1]$ under standard convexity conditions (see Supplementary Note~1 for formal analysis; values $\recoveryratio > 1$ are possible when the corrected operator provides beneficial regularization).
$\recoveryratio = 0$ indicates that calibration yields no benefit (mismatch
is not the bottleneck), while $\recoveryratio = 1$ indicates that
calibration fully closes the mismatch gap.


%% =========================================================================
\subsection*{Agent System Architecture}
\label{sec:methods:agents}

The three gates are implemented as a deterministic scoring pipeline:
an orchestrator maps the user request to a modality from the registry,
then dispatches three parallel scoring modules---one per gate---whose
reports are fused by a bottleneck classifier that identifies the dominant
gate and by a negotiator that enforces cross-gate veto rules
(e.g., rejecting reconstruction when photon starvation coincides with
aggressive compression).  All quantitative decisions are deterministic;
an LLM is used only as an optional fallback for natural-language modality
mapping and is never on the scoring path.  The full agent architecture
and ablation studies are described in a companion paper.

\paragraph{Contract system.}
Inter-agent communication uses 25 Pydantic v2 contract models.  All
contracts inherit from \texttt{StrictBaseModel}, which enforces
\texttt{extra="forbid"} (no unexpected fields),
\texttt{validate\_assignment=True} (mutations re-validated), and a model
validator that rejects NaN and Inf in any float field.  Bounded scores use
\texttt{Field(ge=0.0, le=1.0)}.  Enums are string enums for human-readable
JSON serialization.  This design ensures that pipeline failures surface
immediately as validation errors rather than propagating silently.

\paragraph{YAML registries.}
The system is driven by 9 YAML registries totalling 7,034 lines:
\texttt{modalities.yaml} (modality definitions),
\texttt{graph\_templates.yaml} (OperatorGraph skeletons),
\texttt{photon\_db.yaml} (photon models),
\texttt{mismatch\_db.yaml} (mismatch parameters and correction methods),
\texttt{compression\_db.yaml} (recoverability tables with provenance),
\texttt{solver\_registry.yaml} (solver configurations),
\texttt{primitives.yaml} (primitive operator metadata),
\texttt{dataset\_registry.yaml} (dataset locations and formats), and
\texttt{acceptance\_thresholds.yaml} (pass/fail thresholds per metric).


%% =========================================================================
\subsection*{Correction Algorithms}
\label{sec:methods:correction}

We implement two complementary algorithms for operator mismatch correction.
Crucially, both algorithms operate on the forward operator parameters
$\boldsymbol{\theta}$ rather than the reconstruction solver weights,
making them \emph{solver-agnostic}: the corrected operator
$H(\hat{\boldsymbol{\theta}})$ benefits any downstream solver (GAP-TV,
MST-L, HDNet~\cite{hu2022hdnet}, CST, \etc{}) without retraining.

\paragraph{Algorithm~1: Hierarchical Beam Search.}
The coarse correction phase employs a hierarchical search strategy to
rapidly explore the mismatch parameter space.  For CASSI, the five-parameter
mismatch model comprises mask affine parameters (spatial shifts $dx$, $dy$
and rotation $\theta$) and dispersion parameters (slope $a_1$ and axis
angle $\alpha$); an optional sixth parameter, PSF width $\sigma_{\text{psf}}$, is available but not used in the primary experiments.
The algorithm proceeds as follows:

\begin{enumerate}
  \item \textbf{1D sweeps.}  Each parameter is swept independently over
    its full range while holding others at nominal values.  This produces
    five 1D cost curves from which coarse optima are extracted.
  \item \textbf{3D beam search.}  The mask affine subspace
    $(dx, dy, \theta)$ is searched over a $5 \times 5 \times 5$ grid
    centered on the 1D optima.  The top-$k$ ($k = 5$) candidates by
    reconstruction PSNR are retained.
  \item \textbf{2D beam search.}  For each retained mask candidate,
    the dispersion subspace $(a_1, \alpha)$ is searched over a
    $5 \times 7$ grid.  The joint top-$k$ candidates are retained.
  \item \textbf{Coordinate descent refinement.}  Three rounds of
    univariate refinement on each parameter, shrinking the search
    interval by factor~2 at each round, produce the final estimate
    $\hat{\boldsymbol{\theta}}_{\text{Alg1}}$.
\end{enumerate}

Total runtime is approximately 300 seconds per scene on a single GPU.
Accuracy is $\pm 0.1$--$0.2$ pixels for spatial parameters and
$\pm 0.05^\circ$ for angular parameters.

\paragraph{Algorithm~2: Joint Gradient Refinement.}
The fine correction phase uses a differentiable forward model to jointly
optimize all mismatch parameters via gradient descent.  The key components
are:

\begin{enumerate}
  \item \textbf{Differentiable mask warp.}  The binary mask is warped by a
    continuous affine transformation using bilinear interpolation,
    implemented as a custom PyTorch module
    (\texttt{DifferentiableMaskWarpFixed}).  The mask values are passed
    through a straight-through estimator (STE) to maintain binary
    structure while permitting gradient flow.
  \item \textbf{Differentiable forward model.}  The CASSI forward model
    $y = \texttt{CASSI}(x;\, \boldsymbol{\theta})$ is implemented as a
    differentiable PyTorch module (\texttt{DifferentiableCassiForwardSTE})
    that accepts mismatch parameters as differentiable inputs.
  \item \textbf{GPU grid initialization.}  A full-range 3D grid search
    over $(dx, dy, \theta)$ with $9 \times 9 \times 7 = 567$ points
    provides diverse starting candidates.  The top 9 candidates seed
    multi-start gradient refinement.
  \item \textbf{Staged gradient refinement.}  Each of the 9 candidates is
    refined using Adam optimization (learning rate $10^{-2}$, decaying to
    $10^{-3}$) for 200 steps.  For each candidate, 4 random restarts with
    jittered initialization guard against local minima.  The loss function
    is the negative PSNR computed via an unrolled $K$-iteration
    differentiable GAP-TV solver (\texttt{DifferentiableGAPTV}, $K=10$
    unrolled iterations).
\end{enumerate}

Total runtime for Algorithm~2 is approximately 3,200 seconds
(200~steps $\times$ 4~restarts $\times$ 9~candidates with early stopping).
Accuracy improves to $\pm 0.05$--$0.1$ pixels, a 3--5$\times$ improvement
over Algorithm~1.  The two algorithms are used sequentially in practice:
Algorithm~1 provides a warm start, and Algorithm~2 refines to sub-pixel
precision.


%% =========================================================================
\subsection*{Evaluation Protocol}
\label{sec:methods:eval}

\paragraph{Four-Scenario Protocol.}
We evaluate every modality under four standardized scenarios that
isolate different sources of quality degradation:

\begin{description}
  \item[Scenario~I (Ideal):]
    $\yobs = \Htrue\, \xgt$; reconstruct with $\Htrue$.
    In this scenario the system is perfectly calibrated ($\Htrue = \Hnom$),
    so the operator used for reconstruction matches the one that generated
    the data.  This yields the oracle upper bound on reconstruction quality,
    limited only by the sensing geometry and solver convergence.
  \item[Scenario~II (Mismatch):]
    $\yobs = \Htrue\, \xgt$; reconstruct with $\Hnom$
    ($\Hnom \neq \Htrue$).
    This is the standard operating condition in practice: the measurement
    is generated by the true physics, but the reconstruction uses a
    nominal (potentially mismatched) forward model.
  \item[Scenario~III (Oracle):]
    Same measurements as Scenario~II
    ($\yobs = \Htrue\, \xgt$ with $\Htrue \neq \Hnom$);
    reconstruct with $\Htrue$ instead of $\Hnom$.
    Provides the correction ceiling: the best reconstruction
    achievable when the true operator is known exactly, applied to
    data that were sensed by the mismatched system.  The gap between
    Scenario~III and Scenario~I reveals the irreducible loss from
    the degraded sensing configuration itself (e.g., a shifted mask
    pattern is suboptimal even when perfectly known).
  \item[Scenario~IV (Corrected):]
    $\yobs = \Htrue\, \xgt$; reconstruct with
    $\hat{H} = H(\hat{\boldsymbol{\theta}})$ where
    $\hat{\boldsymbol{\theta}}$ is estimated by Algorithms~1 and~2.
    This quantifies the benefit of mismatch calibration.
\end{description}

\paragraph{Metrics.}
Reconstruction quality is assessed using three complementary metrics:
\begin{itemize}
  \item \textbf{PSNR} (peak signal-to-noise ratio, in dB): the primary
    metric, computed per scene and averaged.  For signals normalized to
    $[0, 1]$, $\psnr = 10 \log_{10}(1 / \text{MSE})$.  For SPC data
    normalized to $[0, 255]$, the peak value is 255.
  \item \textbf{SSIM} (structural similarity index): captures perceptual
    quality including luminance, contrast, and structural components,
    computed with a Gaussian window of width 11 and standard deviation
    1.5.
  \item \textbf{SAM} (spectral angle mapper): for hyperspectral
    modalities (CASSI), measures the angle between predicted and true
    spectral vectors at each spatial location, reported in degrees.
    Lower is better.
\end{itemize}

\paragraph{Datasets.}
\begin{itemize}
  \item \textbf{CASSI:} 10 scenes from the KAIST dataset~\cite{choi2017kaist},
    each a $256 \times 256 \times 28$ spectral cube (28 spectral bands
    from 450\,nm to 650\,nm).  Data range $[0, 1]$.
  \item \textbf{CACTI:} 6 benchmark videos, each $256 \times 256 \times 8$
    (8 temporal frames encoded per snapshot).  Data range $[0, 1]$.
  \item \textbf{SPC:} 11 natural images from the Set11 benchmark, each
    $256 \times 256$ grayscale.  Data range $[0, 255]$.
\end{itemize}

\noindent All per-scene metrics are reported individually as well as
averaged, and all reconstruction arrays are saved as NumPy NPZ files.


%% =========================================================================
\subsection*{Experimental Details}
\label{sec:methods:experimental}

\paragraph{Hardware.}
All experiments are conducted on a single NVIDIA GPU.  Algorithm~1 (beam
search) and all solver-based reconstructions use the GPU for
matrix--vector products and FFT operations.  Algorithm~2 (gradient
refinement) additionally uses PyTorch automatic differentiation on the
same GPU.

\paragraph{CASSI configuration.}
The coded aperture snapshot spectral imaging (CASSI) system uses a TSA-Net
binary mask of dimensions $256 \times 256$, with 28 spectral bands
dispersed along the spatial dimension.  The five-parameter mismatch model
$\mismatchparam = (dx, dy, \theta, a_1, \alpha)$
describes: mask spatial shift in $x$ ($dx$, pixels), mask spatial shift in
$y$ ($dy$, pixels), mask rotation angle ($\theta$, degrees), dispersion
slope ($a_1$, pixels per band), and dispersion axis angle ($\alpha$,
degrees).  An optional sixth parameter, PSF blur width ($\sigma_{\text{psf}}$, pixels), is available but not used in the primary experiments.
These mismatch parameter values were determined through systematic characterization of realistic CASSI assembly errors (Supplementary Note~8). The true mismatch parameters are
$\mismatchparam_{\text{true}} = (dx = 0.5\;\text{px},\; dy = 0.3\;\text{px},\;
\theta = 0.1^\circ,\; a_1 = 2.02,\;
\alpha = 0.15^\circ)$.
Solvers evaluated include TwIST~\cite{bioucasdias2007twist}, GAP-TV~\cite{yuan2016gaptv}, DGSMP~\cite{huang2021dgsmp}, MST-L~\cite{cai2022mst}, and CST-L~\cite{cai2022cst}, all of which receive the same operator and differ only in their reconstruction algorithm. The supplementary per-scene analysis additionally includes DeSCI~\cite{liu2019rank} and HDNet~\cite{hu2022hdnet}.
\textit{Real-data validation.}
The TSA real hyperspectral dataset~\cite{wagadarikar2008cassi} consists of
5 scenes at $660 \times 660$ spatial resolution with 28 spectral bands
and mask-shift step~2.  Four solvers are evaluated: GAP-TV (200 iterations),
HDNet (pre-trained checkpoint, full spatial resolution), MST-S and MST-L
(pre-trained checkpoints, centre-cropped to $256 \times 256$ due to
hardcoded spatial assumptions in the model architecture).  The coded
aperture mask is perturbed by $dx = 0.5$\,px, $dy = 0.3$\,px to simulate
assembly-induced mismatch.  No ground truth is available; quality is
assessed via the normalised measurement residual
$r = \|\mathbf{y} - H\hat{\mathbf{x}}\|^2 / \|\mathbf{y}\|^2$.

\paragraph{CACTI configuration.}
The coded aperture compressive temporal imaging system uses binary
temporal masks of dimensions $256 \times 256$, encoding 8 video frames
into a single snapshot measurement.  Mismatch is parameterized as a
temporal mask timing offset (sub-frame shift).  The default solver is
EfficientSCI~\cite{wang2023efficientsci}.
\textit{Real-data validation.}
The EfficientSCI real temporal dataset~\cite{wang2023efficientsci} consists
of 4 dynamic scenes (duomino, hand, pendulumBall, waterBalloon) at
$512 \times 512$ with compression ratio~10.  The real mask is stored
separately from the measurement data.  Two solvers are evaluated:
GAP-TV (50 iterations) and PnP-FFDNet (50 iterations with FFDNet denoiser).
Mismatch is induced by shifting the mask by $dx = 0.5$\,px, $dy = 0.3$\,px.
Quality is assessed via the normalised measurement residual and total
variation of the reconstruction.

\paragraph{SPC configuration.}
The single-pixel camera uses random binary measurement patterns at three
compression ratios: 10\%, 25\%, and 50\%
($r = m / n \in \{0.10, 0.25, 0.50\}$).  Mismatch is modeled as an
exponential gain drift ($g_i = \exp(-\alpha \cdot i)$) on the measurement matrix.  The default solver
is FISTA-TV with total-variation regularization.

\paragraph{Lensless imaging configuration.}
Lensless cameras replace conventional optics with a diffuser placed directly
on the sensor, making the forward model a 2D convolution with the measured
point-spread function (PSF).  The object is a $64 \times 64$ image
convolved with a Gaussian PSF of width $\sigma = 2.0$\,px.  Mismatch is
parameterized as PSF width error
$\Delta\sigma \in \{0.5, 1.0, 2.0, 3.0\}$\,px, simulating an error in
the assumed pinhole-to-sensor distance.  Five structurally diverse phantoms
are used ($N{=}5$).  The default solver is FISTA-TV (80 iterations).
Calibration uses forward-model fitting on a known reference target:
$\sigma_{\mathrm{cal}} = \mathrm{argmin}_\sigma \|y_{\mathrm{cal}} -
h_\sigma * x_{\mathrm{cal}}\|^2$.
\textit{Real-data validation.}
Five DiffuserCam DLMD Mirflickr scenes~\cite{Antipa2018} are used to
corroborate simulation results on physical photographic content.

\paragraph{Compressive holography configuration.}
Multi-depth inline holography encodes a $(4 \times 64 \times 64)$ 3D
object into a single $(64 \times 64)$ hologram via Fresnel propagation at
four depth planes with spacing $200$\,\textmu m, wavelength
$\lambda = 532$\,nm, and pixel pitch $5$\,\textmu m.  Mismatch is
parameterized as propagation distance error
$\Delta z \in \{10, 50, 100, 200\}$\,\textmu m applied uniformly to all
depth planes.  The default solver is FISTA-TV (80 iterations,
$\lambda_{\mathrm{TV}} = 0.005$).  Calibration uses hologram residual
minimisation: grid search over distance error candidates with the search
range clamped to ensure all propagation distances remain positive.

\paragraph{Fluorescence microscopy configuration.}
Widefield fluorescence imaging uses a dual-PSF Stokes-shift model with
Gaussian excitation PSF ($\sigma_{\mathrm{ex}} = 1.5$\,px) and emission
PSF ($\sigma_{\mathrm{em}} = 2.0$\,px), quantum yield $\eta = 0.7$, and
background level $b = 0.02$.  The specimen is a $64 \times 64$ image with
fluorescent puncta, diffuse organelles, and filamentary structures.
Mismatch is parameterized as PSF sigma error
$\Delta\sigma \in \{0.1, 0.3, 0.5, 1.0\}$\,px applied to both excitation
and emission PSFs.  The default solver is Richardson--Lucy (80 iterations).
Calibration uses a $9 \times 9$ grid search over
$(\sigma_{\mathrm{ex}}, \sigma_{\mathrm{em}})$ centered at the nominal
values, selecting the pair that minimizes the measurement residual.

\paragraph{CT configuration.}
Fan-beam geometry with 180 projections over $180^\circ$.  Mismatch is
modeled as a center-of-rotation (CoR) offset, which produces
characteristic arc artifacts in the reconstruction.  The default solver
is filtered back-projection (FBP)~\cite{feldkamp1984practical} with a Ram-Lak filter, supplemented
by iterative SART for comparison.

\paragraph{CT detector offset configuration.}
Parallel-beam CT is simulated with 360 projection angles over $[0, \pi)$
and 128 detector elements.  The object is a $128 \times 128$ Shepp--Logan
phantom.  Detector offset mismatch is simulated by shifting the sinogram
along the detector axis by $\Delta s \in \{2, 5, 10, 20\}$\,pixels.  The
default solver is filtered backprojection (FBP) with Shepp--Logan filter.
Calibration uses grid search over 21 offset candidates in $[-25, 25]$\,px,
selecting the value that maximizes reconstruction PSNR.

\paragraph{Ptychography configuration.}
Electron ptychography uses 4D-STEM data: a focused probe is scanned across
a $128 \times 128$ grid, recording a convergent-beam electron diffraction
(CBED) pattern at each position (300\,kV, SrTiO$_3$ [001] from
Zenodo~5113449).  The forward model is Fresnel near-field propagation of
the exit wave through the specimen.  Mismatch is parameterized as probe
position jitter $\Delta p \in \{0.5, 1.0, 2.0, 4.0\}$\,px, simulating
scan-coil hysteresis or sample drift.  The default solver is integrated
centre-of-mass (iCoM) phase retrieval.  Calibration uses ePIE blind
position correction~\cite{maiden2009ptychography}.

\paragraph{Cryo-EM configuration.}
Cryo-EM imaging is modeled as contrast transfer function (CTF) filtering
of a 2D projected potential ($128 \times 128$, pixel size $0.1$\,nm).  The
CTF is parameterized by defocus $\Delta f = -2000$\,nm, spherical aberration
$C_s = 2.0$\,mm, electron wavelength $\lambda = 0.00251$\,nm (300\,kV),
with a B-factor envelope ($B = 2$\,nm$^2$) and ice-thickness attenuation
(mean free path $350$\,nm, thickness $50$\,nm).  Mismatch is parameterized
as defocus error $\Delta(\Delta f) \in \{100, 200, 500, 1000\}$\,nm.  The
default solver is Wiener filtering with SNR\,$= 50$.  Calibration uses grid
search over 51 defocus values in $[-3000, -500]$\,nm, selecting the value
that maximizes reconstruction PSNR.

\paragraph{MRI configuration.}
Cartesian $k$-space sampling with 4$\times$ acceleration (25\% of
$k$-space lines acquired).  Mismatch is parameterized as a 5\% multiplicative
error in the coil sensitivity maps used for parallel imaging reconstruction.
The default solver is CG-SENSE~\cite{pruessmann1999sense} with $\ell_1$-wavelet regularization ($\lambda = 10^{-3}$, 30 iterations).
\textit{Scenario~IV calibration (ESPIRiT-adopted).}
For MRI, Scenario~IV adopts ESPIRiT~\cite{uecker2014espirit} as the calibration method: coil sensitivity maps are estimated directly from the central $k$-space autocalibration signal (ACS) region via eigenvalue decomposition of the calibration matrix, then used in the CG-SENSE reconstruction.  When ACS data is sufficient ($\geq 48$ lines), ESPIRiT achieves 85--95\% of the oracle correction ceiling.  For data-limited regimes ($< 32$ ACS lines) where ESPIRiT degrades, the fallback is a grid search over per-coil amplitude and phase scaling (beam-search); this fallback is also the Sc.~IV method for all modalities without an established specialist calibration algorithm (CASSI, CACTI, compressive holography, ultrasound).  The protocol evaluates three conditions---Sc.~I (true maps), Sc.~II (5\% mismatched maps), and Sc.~IV (ESPIRiT-calibrated maps)---on the same CG-SENSE solver to isolate calibration quality from solver design.

\paragraph{Ultrasound configuration.}
Pulse-echo B-mode imaging is simulated with a 64-element linear array,
5\,MHz center frequency, element pitch $0.3$\,mm, sampling rate 100\,MHz
with 1024 time samples, and speed of sound (SoS) $c = 1540$\,m/s.  The
tissue phantom is a $128 \times 128$ reflectivity map with
Rayleigh-distributed speckle, point scatterers, and anechoic cyst regions.
Frequency-dependent attenuation follows
$\alpha(f) = 0.5\,\text{dB\,cm}^{-1}\text{MHz}^{-1}$.  Mismatch is
parameterized as SoS error $\Delta c \in \{50, 100, 200, 400\}$\,m/s,
perturbing time-of-flight calculations.  The default solver is delay-and-sum
(DAS) beamforming (operator adjoint).  Calibration uses grid search over
21 SoS candidates in $[1440, 1640]$\,m/s, selecting the value that
maximizes reconstruction PSNR.

\paragraph{Controlled hardware experiment protocol.}
The software-perturbation protocol above applies calibrated mask shifts to
existing real measurements.  A full hardware-in-the-loop validation
requires physically displacing the coded aperture mask and re-acquiring
data.  The protocol proceeds as follows: (i)~acquire a baseline dataset
with the mask at its factory-calibrated position; (ii)~physically translate
the mask by a known displacement ($\Delta x \in \{0.25, 0.5, 1.0\}$\,px
equivalent, verified by micrometer stage) and re-acquire under identical
illumination; (iii)~reconstruct both datasets with the factory mask
specification and compute the PSNR degradation and measurement residual;
(iv)~apply \pwm{} autonomous calibration and measure recovery.  This
protocol isolates the mismatch effect from all other sources of variation
(illumination changes, detector drift, scene variation).  Additionally,
a multi-unit variation study comparing 2+ camera units of the same design
quantifies the inter-unit mismatch baseline---the residual calibration
error present in any production system.  Clinical CT phantom validation
(ACR Gammex 464) and clinical MRI validation (fastMRI~\cite{zbontar2018fastmri})
configurations are detailed in Supplementary Note~S25.

\paragraph{Real experimental data.}
Three of the five multi-phantom modalities are validated on real
experimental data from established benchmarks: ultrasound uses PICMUS
experimental phantom recordings and in-vivo carotid
scans~\cite{liebgott2016picmus} plus one DeepUS CIRS-040GSE
acquisition~\cite{khan2020deepus}; cryo-EM uses five EMDB 3D density maps
(EMD-5778, EMD-2984, EMD-6287, EMD-11103, EMD-21375) projected at random
orientations; CT uses three LoDoPaB-CT patient
slices~\cite{leuschner2021lodopab}, one FIPS walnut \textmu CT
scan~\cite{hamalainen2015fips}, and one HTC\,2022
sample~\cite{meaney2022htc}.  Compressive holography and fluorescence
microscopy retain synthetic multi-phantom benchmarks as no public datasets
with calibrated ground truth are available for these modalities.

%% =========================================================================
\subsection*{Statistical Analysis}
\label{sec:methods:statistics}

\paragraph{Per-scene reporting.}
All metrics are reported per scene, not merely as dataset averages.  This
enables identification of scene-dependent failure modes (\eg{}, spectrally
flat scenes that are inherently harder for CASSI, or textureless regions
that challenge SPC).

\paragraph{Summary statistics.}
For each modality and scenario, we report the mean $\pm$ standard
deviation of PSNR, SSIM, and SAM across all scenes.  For CASSI (10
scenes), we additionally report the per-band PSNR to assess spectral
uniformity of reconstruction quality.

\paragraph{Recovery ratio confidence intervals.}
The recovery ratio $\recoveryratio$ (\cref{eq:recovery_ratio}) is a ratio
of differences and therefore sensitive to noise in the constituent PSNR
values.  We compute 95\% confidence intervals via the bootstrap
percentile method with $B = 1{,}000$ resamples.  At each bootstrap
iteration, we resample the scene set with replacement, recompute the
mean PSNR for each scenario, and derive $\recoveryratio$.  The 2.5th
and 97.5th percentiles of the bootstrap distribution define the 95\% CI.

\paragraph{Parameter recovery accuracy.}
For mismatch correction experiments, we report the root-mean-square error
(RMSE) between the estimated and true mismatch parameters:
\begin{equation}
  \text{RMSE}_k = \sqrt{\frac{1}{N_{\text{scene}}}
    \sum_{i=1}^{N_{\text{scene}}}
    (\hat\theta_{k,i} - \theta_{k,\text{true}})^2}
\end{equation}
where $k$ indexes the mismatch parameter, $i$ indexes the scene, and
$N_{\text{scene}}$ is the number of test scenes.  Uncertainty in the RMSE
is estimated via bootstrap ($B = 1{,}000$).

\paragraph{Ablation significance.}
Ablation studies (removal of the photon-budget, recoverability, or
mismatch scoring module, or RunBundle discipline) are evaluated by comparing the
full-pipeline PSNR against each ablated variant.  We report the PSNR
difference $\Delta\psnr$ per modality and verify that each component
contributes $\geq 0.5$\,dB across all validated modalities, establishing
practical significance.


%% =========================================================================
\subsection*{Code and Data Availability}
\label{sec:methods:availability}

\paragraph{Source code.}
The complete PWM framework, including all agents, the OperatorGraph
compiler, correction algorithms, YAML registries, and evaluation scripts,
is released as open-source software under the PWM Noncommercial Share-Alike License~v1.0 at
\url{https://github.com/integritynoble/Physics_World_Model}.  The
codebase is organized into two Python packages:
\texttt{pwm\_core} (core framework, agents, graph compiler, calibration
algorithms) and \texttt{pwm\_AI\_Scientist} (automated experiment
generation and analysis).

\paragraph{Reconstruction data.}
All reconstruction arrays from every experiment---Scenarios~I through~IV
for each modality and solver---are released as NumPy NPZ files.
Files are stored using Git LFS and require \texttt{allow\_pickle=True}
for loading.  Data ranges are standardized: CASSI and CACTI
reconstructions are normalized to $[0, 1]$; SPC reconstructions are in
$[0, 255]$.

\paragraph{Experiment manifests.}
Every experiment is recorded in a RunBundle v0.3.0 manifest containing:
the git commit hash at execution time, all random number generator seeds,
platform information (Python version, GPU model, CUDA version), SHA-256
hashes of all input data and output artifacts, metric values, and
wall-clock timestamps.  These manifests enable exact reproduction of every
reported result.

\paragraph{Registry data.}
All 9 YAML registries (7,034 lines total) that drive the agent system---including
modality definitions, graph templates, photon models, mismatch databases,
compression tables, solver configurations, primitive specifications,
dataset paths, and acceptance thresholds---are publicly available in the
repository under \texttt{packages/pwm\_core/contrib/}.
The \texttt{ExperimentSpec} JSON schemas used for pipeline input
validation are included alongside worked examples in
\texttt{examples/}.
 from the main manuscript.

\section*{Online Methods}

%% =========================================================================
\subsection*{OperatorGraph Specification}
\label{sec:methods:operatorgraph}

\paragraph{Formal definition.}
The \opg{} intermediate representation encodes the forward physics of any
computational imaging modality as a directed acyclic graph (DAG)
$\mathcal{G} = (\mathcal{V}, \mathcal{E})$.
Each node $v_i \in \mathcal{V}$ wraps a \emph{primitive operator} and
implements two entry points: $\texttt{forward}(x) \to y$ and
$\texttt{adjoint}(y) \to x$, the latter defined only when the primitive
is linear.  Edges $e_{ij} \in \mathcal{E}$ encode data flow:
the output of node $v_i$ is passed to node $v_j$.
Each node additionally exposes a set of learnable parameters
$\boldsymbol{\theta}_i$ that may be perturbed during mismatch simulation
or optimized during calibration, as well as read-only metadata flags
(\texttt{is\_linear}, \texttt{is\_stochastic}, \texttt{is\_differentiable}).
The graph is stored as a declarative YAML specification
(\texttt{OperatorGraphSpec}) and compiled to an executable
\texttt{GraphOperator} object by the \texttt{GraphCompiler}.

\paragraph{Node types.}
Primitive operators fall into two categories:

\begin{itemize}
  \item \textbf{Linear operators.}  Convolution (\texttt{conv2d}), mask
    modulation (\texttt{mask\_modulate}), sub-pixel shift (\texttt{subpixel\_shift\_2d}),
    Radon transform (\texttt{radon\_fanbeam}),
    Fourier encoding (\texttt{fourier\_encode}), spectral dispersion
    (\texttt{spectral\_disperse}), Fresnel propagation
    (\texttt{fresnel\_propagate}), random projection
    (\texttt{random\_project}), and structured illumination
    (\texttt{sim\_modulate}).
    Each implements both \texttt{forward()} and \texttt{adjoint()}.
  \item \textbf{Nonlinear operators.}  Squared magnitude
    (\texttt{magnitude\_sq}), Poisson--Gaussian noise
    (\texttt{poisson\_gaussian}), saturation clipping
    (\texttt{saturation\_clip}), phase retrieval nonlinearity
    (\texttt{phase\_abs}), and detector quantization
    (\texttt{quantize}).
    These set \texttt{is\_linear = False} and raise
    \texttt{NotImplementedError} on \texttt{adjoint()}, except where a
    well-defined pseudo-adjoint exists (\eg{}, the identity adjoint for
    magnitude-squared in Gerchberg--Saxton-type algorithms).
\end{itemize}

\paragraph{Adjoint consistency and Lipschitz bounds as mathematical safety rails.}
Two structural constraints distinguish the \opg{} IR from standard deep-learning
forward models and prevent either nonlinear primitive family from becoming a
universal approximator.

First, \emph{adjoint consistency}: correctness of every linear primitive
is verified by a randomized dot-product test.  For a primitive $A$ with
forward map $A:\mathbb{R}^n \to \mathbb{R}^m$, we draw $x \sim \mathcal{N}(0,I_n)$
and $y \sim \mathcal{N}(0,I_m)$ and compute
\begin{equation}
  \delta = \frac{|\langle A^* y,\, x \rangle - \langle y,\, Ax \rangle|}
               {\max(|\langle A^* y,\, x \rangle|,\; \epsilon)}
  \label{eq:adjoint_check}
\end{equation}
where $\epsilon = 10^{-12}$ guards against division by zero.  The test
is repeated $n_{\text{trials}} = 5$ times with independent random draws;
the primitive passes if $\delta_{\max} < 10^{-6}$.  At the graph level,
a compiled \texttt{GraphOperator} composed entirely of linear nodes
executes the same test over the composed forward--adjoint chain.  A
\texttt{GraphAdjointCheckReport} records $n_{\text{trials}}$,
$\delta_{\max}$, and $\bar\delta$ for audit.  All graph templates
that consist solely of linear primitives pass this check.
Adjoint consistency is not merely a numerical correctness guarantee: it
is a \emph{physical faithfulness certificate}.  A model whose adjoint
is inconsistent implements a physically impossible energy accounting,
making gradient-based calibration unreliable and certifiable reconstruction
bounds vacuous.  Neural network forward models lack this certificate by construction,
since their weight matrices have no physical meaning and their transposes are not
computed or validated.

Second, \emph{Lipschitz bounds}: the two nonlinear primitive families,
Detect~$D$ and Transform~$\Lambda$, are each restricted to five canonical
function families with at most two scalar parameters.  This restriction
enforces bounded Lipschitz constants ($\mathrm{Lip}(\Lambda_k) \leq L$ for
the explicit $L$ derived in~\cite{yang2026fpt}) and prevents either primitive
from becoming a universal approximator in the sense of Cybenko~\cite{cybenko1989approximation}.
Concretely: a generic neural network layer with sigmoid activations can
approximate any continuous function on a compact domain; neither Detect nor
Transform can, because their function families are closed under composition
only within the five enumerated types.  This non-universality is the key
property that makes the 11-primitive library \emph{falsifiable}: if a
modality's forward physics cannot be represented within the restricted
families, the extension protocol (Main Text, ``Extension protocol'') fires
and a new primitive must be justified physically.  The combination of
adjoint consistency and bounded Lipschitz nonlinearity provides the
mathematical safety rails that make \opg{}-based reconstruction both
interpretable and certifiable---properties that standard physics-informed
neural networks~\cite{raissi2019pinn} achieve only approximately and
without modality-agnostic guarantees.

\paragraph{Graph compilation.}
The compiler executes a four-stage pipeline:
\begin{enumerate}
  \item \textbf{Validate.}  Confirm acyclicity via topological sort
    (Kahn's algorithm), verify that every \texttt{primitive\_id}
    exists in the global \texttt{PRIMITIVE\_REGISTRY}, reject
    duplicate \texttt{node\_id} values, and optionally verify shape
    compatibility along edges when a \texttt{canonical\_chain}
    metadata flag is set.
  \item \textbf{Bind.}  Instantiate each primitive with its parameter
    dictionary $\boldsymbol{\theta}_i$.
  \item \textbf{Plan forward.}  The topological sort yields a sequential
    execution plan $(v_{\pi(1)}, \ldots, v_{\pi(|\mathcal{V}|)})$.
  \item \textbf{Plan adjoint.}  For graphs where
    \texttt{all\_linear = True}, the adjoint plan reverses the
    topological order and applies each node's individual adjoint in
    sequence, implementing the chain rule $A^* = A_1^*
    \circ \cdots \circ A_{|\mathcal{V}|}^*$ for a composition $A = A_{|\mathcal{V}|}
    \circ \cdots \circ A_1$.  For graphs containing nonlinear nodes,
    the adjoint plan is not generated, and any call to
    \texttt{adjoint()} raises \texttt{NotImplementedError} at runtime.
\end{enumerate}
The compiled \texttt{GraphOperator} is serializable to JSON and hashable
via SHA-256 for provenance tracking in RunBundle manifests.

\paragraph{Template library.}
The \texttt{graph\_templates.yaml} registry contains 170 templates
spanning all five carrier families. Of these, 12 modalities carry full
end-to-end Scenario~I--IV correction validation across all five carrier
families (Supplementary Table~S3). Representative modalities by carrier family include:

\begin{itemize}
  \item \textbf{Optical and X-ray photons:}  CASSI, SPC, CACTI, structured
    illumination microscopy (SIM), confocal, light-sheet, holography,
    ptychography, Fourier ptychographic microscopy (FPM), optical coherence
    tomography (OCT), lensless imaging, light field, integral imaging,
    neural radiance fields (NeRF), Gaussian splatting, fluorescence
    lifetime imaging (FLIM), diffuse optical tomography (DOT),
    phase retrieval, X-ray computed tomography (CT), and cone-beam CT
    (CBCT).
  \item \textbf{Electrons:}  Electron diffraction, electron backscatter
    diffraction (EBSD), electron energy loss spectroscopy (EELS), and
    electron holography.
  \item \textbf{Nuclear spins (MRI):}  Functional MRI (fMRI), diffusion-weighted
    MRI (DW-MRI), and magnetic resonance spectroscopy (MRS).
  \item \textbf{Acoustic waves:}  Ultrasound B-mode, Doppler ultrasound,
    shear-wave elastography, sonar, and photoacoustic tomography
    (combines optical excitation with acoustic detection).
  \item \textbf{Particle-beam probes:}  Neutron tomography, proton CT,
    and muon tomography.  Each decomposes into existing primitives with
    no extension required: neutron CT follows the X-ray CT DAG
    ($P \!\to\! \Lambda \!\to\! D$, Beer--Lambert attenuation with nuclear
    cross-section $\mu_n$); proton CT inserts a Bethe--Bloch
    stopping-power stage as Transform~$\Lambda$ (\#11); muon tomography
    uses Scatter~$R$ (\#10) for multiple Coulomb scattering and
    POCA vertex reconstruction (Supplementary Note~S24).
\end{itemize}

\paragraph{Fidelity levels.}
Each template is parameterized by a fidelity level that controls the
degree of physical realism in the simulated forward model:

\begin{description}
  \item[Level~1 (Linear, shift-invariant):]  The forward model is a
    linear, spatially uniform operator---the simplest approximation,
    suitable for initial diagnostics and rapid prototyping.
  \item[Level~2 (Linear, shift-variant):]  Spatially varying operator
    parameters (e.g.\ non-uniform illumination, position-dependent PSF,
    multi-coil sensitivity maps in MRI).  Adds a modality-appropriate
    noise model (Poisson shot noise plus Gaussian read noise for
    photon-counting modalities, Rician noise for MRI, Poisson for CT).
  \item[Level~3 (Nonlinear, ray/wave-based):]  Includes nonlinear
    effects such as beam hardening, phase wrapping, and scattering,
    captured by the Transform~$\Lambda$ primitive.
    Perturbation families and ranges are specified in
    \texttt{mismatch\_db.yaml}.
  \item[Level~4 (Full-wave / Monte Carlo):]  Complete physical
    simulation including wave-optical propagation, spatially varying
    aberrations, detector nonlinearities, and environmental drift.
    Currently implemented for holography and ptychography.
\end{description}

\noindent All 11 primitives in the library $\Blib$ apply uniformly across every fidelity level; the levels organize simulation complexity, not theorem scope.


%% =========================================================================
\subsection*{Triad Decomposition Formalization}
\label{sec:methods:triad}

The \triad{} asserts that the quality of any computational imaging
reconstruction is bounded by three fundamental gates.  Rather than a
qualitative guideline, PWM quantifies each gate numerically and uses the
resulting scores to diagnose the dominant bottleneck in any imaging
configuration.

\paragraph{Gate~1 (Recoverability).}
Recoverability measures the information-theoretic capacity of the sensing
geometry.  We quantify it via the \emph{effective compression ratio}
$r = m / n$, where $m$ is the number of independent measurements and $n$
the dimension of the signal.  The \texttt{compression\_db.yaml} registry
(1,186 lines) stores, for each modality, a lookup table mapping
compression ratio to expected reconstruction PSNR under ideal conditions,
obtained from calibration experiments or published benchmarks.  Each entry
carries a \texttt{provenance} field citing the source (paper DOI, internal
experiment ID, or theoretical formula).  Additional recoverability
indicators include the effective rank of the measurement matrix (estimated
via randomized SVD for large operators), the dimension of the null space,
and the restricted isometry property (RIP) constant where analytically
tractable (\eg{}, for Gaussian random projections in SPC).

\paragraph{Gate~2 (Carrier Budget).}
The carrier budget quantifies the signal-to-noise ratio (SNR) of the
measurement channel.  The photon-budget module consumes the
\texttt{photon\_db.yaml} registry which stores, per modality,
a deterministic photon model parameterized by source power, quantum
efficiency, exposure time, and detector characteristics.  It
classifies the noise regime into one of three categories:
\emph{shot-limited} (Poisson-dominated, SNR $\propto \sqrt{N_{\text{photon}}}$),
\emph{read-limited} (Gaussian read noise dominates, SNR $\propto
N_{\text{photon}} / \sigma_{\text{read}}$), and \emph{dark-current-limited}
(long exposures where dark current accumulation dominates).  The output is
a \texttt{PhotonReport} containing the estimated SNR in decibels, the noise
regime classification, per-element photon count, and a feasibility verdict
(\texttt{sufficient}, \texttt{marginal}, or \texttt{insufficient}).

\paragraph{Gate~3 (Operator Mismatch).}
Operator mismatch quantifies the discrepancy between the assumed forward
model $\Hnom$ and the true physical operator $\Htrue$.  The
mismatch scoring module consults \texttt{mismatch\_db.yaml}
which catalogs, for each modality, the set of mismatch parameters
(spatial shifts, rotational offsets, dispersion errors, PSF deviations,
coil sensitivity errors, center-of-rotation offsets, \etc{}), their
typical ranges, and available correction methods.  The mismatch severity
score $s \in [0, 1]$ is computed as the normalized $\ell_2$ distance
$\| \boldsymbol{\theta}_{\text{true}} - \boldsymbol{\theta}_{\text{nom}} \|
/ \| \boldsymbol{\theta}_{\text{range}} \|$, where
$\boldsymbol{\theta}_{\text{range}}$ is the per-parameter dynamic range
from the registry.  Sensitivity analysis $\partial \psnr / \partial
\theta_k$ is estimated via finite differences on the forward model.  The
output is a \texttt{MismatchReport} containing the severity score, the
dominant mismatch parameter, the recommended correction method, and the
expected PSNR gain from correction.

\paragraph{Gate binding determination.}
Given reconstruction results under the four-scenario protocol
(the Evaluation Protocol section below), PWM identifies the dominant gate by
comparing three cost terms:
\begin{align}
  C_{\text{mismatch}} &= \psnr_{\text{I}} - \psnr_{\text{II}} \label{eq:cost_mismatch} \\
  C_{\text{noise}}    &= \psnr_{\text{ideal}} - \psnr_{\text{noisy}} \label{eq:cost_noise} \\
  C_{\text{recover}}  &= \psnr_{\text{limit}} - \psnr_{\text{I}} \label{eq:cost_recover}
\end{align}
where $\psnr_{\text{I}}$ is the reconstruction PSNR under Scenario~I
(ideal operator), $\psnr_{\text{II}}$ under Scenario~II (mismatched
operator), $\psnr_{\text{noisy}}$ under the corresponding noisy condition,
and $\psnr_{\text{limit}}$ is the theoretical upper bound from the
compression table.  The dominant gate is $\arg\max_g C_g$.

\paragraph{TriadReport.}
For every diagnosis, \pwm{} produces a \triadreport{}: a Pydantic-validated
structured artifact comprising
\texttt{dominant\_gate} (enum: \texttt{recoverability},
\texttt{carrier\_budget}, \texttt{operator\_mismatch}),
\texttt{evidence\_scores} (three floats, one per gate),
\texttt{confidence\_interval} (float, 95\% CI width from bootstrap),
\texttt{recommended\_action} (string, \eg{} ``increase compression ratio''
or ``apply mismatch correction''), and
\texttt{parameter\_sensitivities} (dictionary mapping each mismatch
parameter name to its $\partial \psnr / \partial \theta_k$ value).
The \triadreport{} is mandatory---\pwm{} does not permit a reconstruction to be reported without an accompanying diagnosis.

\paragraph{Recovery ratio.}
We define the \emph{recovery ratio}
\begin{equation}
  \recoveryratio = \frac{\psnr_{\text{IV}} - \psnr_{\text{II}}}
                        {\psnr_{\text{I}} - \psnr_{\text{II}}}
  \label{eq:recovery_ratio}
\end{equation}
which lies in $[0, 1]$ under standard convexity conditions (see Supplementary Note~1 for formal analysis; values $\recoveryratio > 1$ are possible when the corrected operator provides beneficial regularization).
$\recoveryratio = 0$ indicates that calibration yields no benefit (mismatch
is not the bottleneck), while $\recoveryratio = 1$ indicates that
calibration fully closes the mismatch gap.


%% =========================================================================
\subsection*{Agent System Architecture}
\label{sec:methods:agents}

The three gates are implemented as a deterministic scoring pipeline:
an orchestrator maps the user request to a modality from the registry,
then dispatches three parallel scoring modules---one per gate---whose
reports are fused by a bottleneck classifier that identifies the dominant
gate and by a negotiator that enforces cross-gate veto rules
(e.g., rejecting reconstruction when photon starvation coincides with
aggressive compression).  All quantitative decisions are deterministic;
an LLM is used only as an optional fallback for natural-language modality
mapping and is never on the scoring path.  The full agent architecture
and ablation studies are described in a companion paper.

\paragraph{Contract system.}
Inter-agent communication uses 25 Pydantic v2 contract models.  All
contracts inherit from \texttt{StrictBaseModel}, which enforces
\texttt{extra="forbid"} (no unexpected fields),
\texttt{validate\_assignment=True} (mutations re-validated), and a model
validator that rejects NaN and Inf in any float field.  Bounded scores use
\texttt{Field(ge=0.0, le=1.0)}.  Enums are string enums for human-readable
JSON serialization.  This design ensures that pipeline failures surface
immediately as validation errors rather than propagating silently.

\paragraph{YAML registries.}
The system is driven by 9 YAML registries totalling 7,034 lines:
\texttt{modalities.yaml} (modality definitions),
\texttt{graph\_templates.yaml} (OperatorGraph skeletons),
\texttt{photon\_db.yaml} (photon models),
\texttt{mismatch\_db.yaml} (mismatch parameters and correction methods),
\texttt{compression\_db.yaml} (recoverability tables with provenance),
\texttt{solver\_registry.yaml} (solver configurations),
\texttt{primitives.yaml} (primitive operator metadata),
\texttt{dataset\_registry.yaml} (dataset locations and formats), and
\texttt{acceptance\_thresholds.yaml} (pass/fail thresholds per metric).


%% =========================================================================
\subsection*{Correction Algorithms}
\label{sec:methods:correction}

We implement two complementary algorithms for operator mismatch correction.
Crucially, both algorithms operate on the forward operator parameters
$\boldsymbol{\theta}$ rather than the reconstruction solver weights,
making them \emph{solver-agnostic}: the corrected operator
$H(\hat{\boldsymbol{\theta}})$ benefits any downstream solver (GAP-TV,
MST-L, HDNet~\cite{hu2022hdnet}, CST, \etc{}) without retraining.

\paragraph{Algorithm~1: Hierarchical Beam Search.}
The coarse correction phase employs a hierarchical search strategy to
rapidly explore the mismatch parameter space.  For CASSI, the five-parameter
mismatch model comprises mask affine parameters (spatial shifts $dx$, $dy$
and rotation $\theta$) and dispersion parameters (slope $a_1$ and axis
angle $\alpha$); an optional sixth parameter, PSF width $\sigma_{\text{psf}}$, is available but not used in the primary experiments.
The algorithm proceeds as follows:

\begin{enumerate}
  \item \textbf{1D sweeps.}  Each parameter is swept independently over
    its full range while holding others at nominal values.  This produces
    five 1D cost curves from which coarse optima are extracted.
  \item \textbf{3D beam search.}  The mask affine subspace
    $(dx, dy, \theta)$ is searched over a $5 \times 5 \times 5$ grid
    centered on the 1D optima.  The top-$k$ ($k = 5$) candidates by
    reconstruction PSNR are retained.
  \item \textbf{2D beam search.}  For each retained mask candidate,
    the dispersion subspace $(a_1, \alpha)$ is searched over a
    $5 \times 7$ grid.  The joint top-$k$ candidates are retained.
  \item \textbf{Coordinate descent refinement.}  Three rounds of
    univariate refinement on each parameter, shrinking the search
    interval by factor~2 at each round, produce the final estimate
    $\hat{\boldsymbol{\theta}}_{\text{Alg1}}$.
\end{enumerate}

Total runtime is approximately 300 seconds per scene on a single GPU.
Accuracy is $\pm 0.1$--$0.2$ pixels for spatial parameters and
$\pm 0.05^\circ$ for angular parameters.

\paragraph{Algorithm~2: Joint Gradient Refinement.}
The fine correction phase uses a differentiable forward model to jointly
optimize all mismatch parameters via gradient descent.  The key components
are:

\begin{enumerate}
  \item \textbf{Differentiable mask warp.}  The binary mask is warped by a
    continuous affine transformation using bilinear interpolation,
    implemented as a custom PyTorch module
    (\texttt{DifferentiableMaskWarpFixed}).  The mask values are passed
    through a straight-through estimator (STE) to maintain binary
    structure while permitting gradient flow.
  \item \textbf{Differentiable forward model.}  The CASSI forward model
    $y = \texttt{CASSI}(x;\, \boldsymbol{\theta})$ is implemented as a
    differentiable PyTorch module (\texttt{DifferentiableCassiForwardSTE})
    that accepts mismatch parameters as differentiable inputs.
  \item \textbf{GPU grid initialization.}  A full-range 3D grid search
    over $(dx, dy, \theta)$ with $9 \times 9 \times 7 = 567$ points
    provides diverse starting candidates.  The top 9 candidates seed
    multi-start gradient refinement.
  \item \textbf{Staged gradient refinement.}  Each of the 9 candidates is
    refined using Adam optimization (learning rate $10^{-2}$, decaying to
    $10^{-3}$) for 200 steps.  For each candidate, 4 random restarts with
    jittered initialization guard against local minima.  The loss function
    is the negative PSNR computed via an unrolled $K$-iteration
    differentiable GAP-TV solver (\texttt{DifferentiableGAPTV}, $K=10$
    unrolled iterations).
\end{enumerate}

Total runtime for Algorithm~2 is approximately 3,200 seconds
(200~steps $\times$ 4~restarts $\times$ 9~candidates with early stopping).
Accuracy improves to $\pm 0.05$--$0.1$ pixels, a 3--5$\times$ improvement
over Algorithm~1.  The two algorithms are used sequentially in practice:
Algorithm~1 provides a warm start, and Algorithm~2 refines to sub-pixel
precision.


%% =========================================================================
\subsection*{Evaluation Protocol}
\label{sec:methods:eval}

\paragraph{Four-Scenario Protocol.}
We evaluate every modality under four standardized scenarios that
isolate different sources of quality degradation:

\begin{description}
  \item[Scenario~I (Ideal):]
    $\yobs = \Htrue\, \xgt$; reconstruct with $\Htrue$.
    In this scenario the system is perfectly calibrated ($\Htrue = \Hnom$),
    so the operator used for reconstruction matches the one that generated
    the data.  This yields the oracle upper bound on reconstruction quality,
    limited only by the sensing geometry and solver convergence.
  \item[Scenario~II (Mismatch):]
    $\yobs = \Htrue\, \xgt$; reconstruct with $\Hnom$
    ($\Hnom \neq \Htrue$).
    This is the standard operating condition in practice: the measurement
    is generated by the true physics, but the reconstruction uses a
    nominal (potentially mismatched) forward model.
  \item[Scenario~III (Oracle):]
    Same measurements as Scenario~II
    ($\yobs = \Htrue\, \xgt$ with $\Htrue \neq \Hnom$);
    reconstruct with $\Htrue$ instead of $\Hnom$.
    Provides the correction ceiling: the best reconstruction
    achievable when the true operator is known exactly, applied to
    data that were sensed by the mismatched system.  The gap between
    Scenario~III and Scenario~I reveals the irreducible loss from
    the degraded sensing configuration itself (e.g., a shifted mask
    pattern is suboptimal even when perfectly known).
  \item[Scenario~IV (Corrected):]
    $\yobs = \Htrue\, \xgt$; reconstruct with
    $\hat{H} = H(\hat{\boldsymbol{\theta}})$ where
    $\hat{\boldsymbol{\theta}}$ is estimated by Algorithms~1 and~2.
    This quantifies the benefit of mismatch calibration.
\end{description}

\paragraph{Metrics.}
Reconstruction quality is assessed using three complementary metrics:
\begin{itemize}
  \item \textbf{PSNR} (peak signal-to-noise ratio, in dB): the primary
    metric, computed per scene and averaged.  For signals normalized to
    $[0, 1]$, $\psnr = 10 \log_{10}(1 / \text{MSE})$.  For SPC data
    normalized to $[0, 255]$, the peak value is 255.
  \item \textbf{SSIM} (structural similarity index): captures perceptual
    quality including luminance, contrast, and structural components,
    computed with a Gaussian window of width 11 and standard deviation
    1.5.
  \item \textbf{SAM} (spectral angle mapper): for hyperspectral
    modalities (CASSI), measures the angle between predicted and true
    spectral vectors at each spatial location, reported in degrees.
    Lower is better.
\end{itemize}

\paragraph{Datasets.}
\begin{itemize}
  \item \textbf{CASSI:} 10 scenes from the KAIST dataset~\cite{choi2017kaist},
    each a $256 \times 256 \times 28$ spectral cube (28 spectral bands
    from 450\,nm to 650\,nm).  Data range $[0, 1]$.
  \item \textbf{CACTI:} 6 benchmark videos, each $256 \times 256 \times 8$
    (8 temporal frames encoded per snapshot).  Data range $[0, 1]$.
  \item \textbf{SPC:} 11 natural images from the Set11 benchmark, each
    $256 \times 256$ grayscale.  Data range $[0, 255]$.
\end{itemize}

\noindent All per-scene metrics are reported individually as well as
averaged, and all reconstruction arrays are saved as NumPy NPZ files.


%% =========================================================================
\subsection*{Experimental Details}
\label{sec:methods:experimental}

\paragraph{Hardware.}
All experiments are conducted on a single NVIDIA GPU.  Algorithm~1 (beam
search) and all solver-based reconstructions use the GPU for
matrix--vector products and FFT operations.  Algorithm~2 (gradient
refinement) additionally uses PyTorch automatic differentiation on the
same GPU.

\paragraph{CASSI configuration.}
The coded aperture snapshot spectral imaging (CASSI) system uses a TSA-Net
binary mask of dimensions $256 \times 256$, with 28 spectral bands
dispersed along the spatial dimension.  The five-parameter mismatch model
$\mismatchparam = (dx, dy, \theta, a_1, \alpha)$
describes: mask spatial shift in $x$ ($dx$, pixels), mask spatial shift in
$y$ ($dy$, pixels), mask rotation angle ($\theta$, degrees), dispersion
slope ($a_1$, pixels per band), and dispersion axis angle ($\alpha$,
degrees).  An optional sixth parameter, PSF blur width ($\sigma_{\text{psf}}$, pixels), is available but not used in the primary experiments.
These mismatch parameter values were determined through systematic characterization of realistic CASSI assembly errors (Supplementary Note~8). The true mismatch parameters are
$\mismatchparam_{\text{true}} = (dx = 0.5\;\text{px},\; dy = 0.3\;\text{px},\;
\theta = 0.1^\circ,\; a_1 = 2.02,\;
\alpha = 0.15^\circ)$.
Solvers evaluated include TwIST~\cite{bioucasdias2007twist}, GAP-TV~\cite{yuan2016gaptv}, DGSMP~\cite{huang2021dgsmp}, MST-L~\cite{cai2022mst}, and CST-L~\cite{cai2022cst}, all of which receive the same operator and differ only in their reconstruction algorithm. The supplementary per-scene analysis additionally includes DeSCI~\cite{liu2019rank} and HDNet~\cite{hu2022hdnet}.
\textit{Real-data validation.}
The TSA real hyperspectral dataset~\cite{wagadarikar2008cassi} consists of
5 scenes at $660 \times 660$ spatial resolution with 28 spectral bands
and mask-shift step~2.  Four solvers are evaluated: GAP-TV (200 iterations),
HDNet (pre-trained checkpoint, full spatial resolution), MST-S and MST-L
(pre-trained checkpoints, centre-cropped to $256 \times 256$ due to
hardcoded spatial assumptions in the model architecture).  The coded
aperture mask is perturbed by $dx = 0.5$\,px, $dy = 0.3$\,px to simulate
assembly-induced mismatch.  No ground truth is available; quality is
assessed via the normalised measurement residual
$r = \|\mathbf{y} - H\hat{\mathbf{x}}\|^2 / \|\mathbf{y}\|^2$.

\paragraph{CACTI configuration.}
The coded aperture compressive temporal imaging system uses binary
temporal masks of dimensions $256 \times 256$, encoding 8 video frames
into a single snapshot measurement.  Mismatch is parameterized as a
temporal mask timing offset (sub-frame shift).  The default solver is
EfficientSCI~\cite{wang2023efficientsci}.
\textit{Real-data validation.}
The EfficientSCI real temporal dataset~\cite{wang2023efficientsci} consists
of 4 dynamic scenes (duomino, hand, pendulumBall, waterBalloon) at
$512 \times 512$ with compression ratio~10.  The real mask is stored
separately from the measurement data.  Two solvers are evaluated:
GAP-TV (50 iterations) and PnP-FFDNet (50 iterations with FFDNet denoiser).
Mismatch is induced by shifting the mask by $dx = 0.5$\,px, $dy = 0.3$\,px.
Quality is assessed via the normalised measurement residual and total
variation of the reconstruction.

\paragraph{SPC configuration.}
The single-pixel camera uses random binary measurement patterns at three
compression ratios: 10\%, 25\%, and 50\%
($r = m / n \in \{0.10, 0.25, 0.50\}$).  Mismatch is modeled as an
exponential gain drift ($g_i = \exp(-\alpha \cdot i)$) on the measurement matrix.  The default solver
is FISTA-TV with total-variation regularization.

\paragraph{Lensless imaging configuration.}
Lensless cameras replace conventional optics with a diffuser placed directly
on the sensor, making the forward model a 2D convolution with the measured
point-spread function (PSF).  The object is a $64 \times 64$ image
convolved with a Gaussian PSF of width $\sigma = 2.0$\,px.  Mismatch is
parameterized as PSF width error
$\Delta\sigma \in \{0.5, 1.0, 2.0, 3.0\}$\,px, simulating an error in
the assumed pinhole-to-sensor distance.  Five structurally diverse phantoms
are used ($N{=}5$).  The default solver is FISTA-TV (80 iterations).
Calibration uses forward-model fitting on a known reference target:
$\sigma_{\mathrm{cal}} = \mathrm{argmin}_\sigma \|y_{\mathrm{cal}} -
h_\sigma * x_{\mathrm{cal}}\|^2$.
\textit{Real-data validation.}
Five DiffuserCam DLMD Mirflickr scenes~\cite{Antipa2018} are used to
corroborate simulation results on physical photographic content.

\paragraph{Compressive holography configuration.}
Multi-depth inline holography encodes a $(4 \times 64 \times 64)$ 3D
object into a single $(64 \times 64)$ hologram via Fresnel propagation at
four depth planes with spacing $200$\,\textmu m, wavelength
$\lambda = 532$\,nm, and pixel pitch $5$\,\textmu m.  Mismatch is
parameterized as propagation distance error
$\Delta z \in \{10, 50, 100, 200\}$\,\textmu m applied uniformly to all
depth planes.  The default solver is FISTA-TV (80 iterations,
$\lambda_{\mathrm{TV}} = 0.005$).  Calibration uses hologram residual
minimisation: grid search over distance error candidates with the search
range clamped to ensure all propagation distances remain positive.

\paragraph{Fluorescence microscopy configuration.}
Widefield fluorescence imaging uses a dual-PSF Stokes-shift model with
Gaussian excitation PSF ($\sigma_{\mathrm{ex}} = 1.5$\,px) and emission
PSF ($\sigma_{\mathrm{em}} = 2.0$\,px), quantum yield $\eta = 0.7$, and
background level $b = 0.02$.  The specimen is a $64 \times 64$ image with
fluorescent puncta, diffuse organelles, and filamentary structures.
Mismatch is parameterized as PSF sigma error
$\Delta\sigma \in \{0.1, 0.3, 0.5, 1.0\}$\,px applied to both excitation
and emission PSFs.  The default solver is Richardson--Lucy (80 iterations).
Calibration uses a $9 \times 9$ grid search over
$(\sigma_{\mathrm{ex}}, \sigma_{\mathrm{em}})$ centered at the nominal
values, selecting the pair that minimizes the measurement residual.

\paragraph{CT configuration.}
Fan-beam geometry with 180 projections over $180^\circ$.  Mismatch is
modeled as a center-of-rotation (CoR) offset, which produces
characteristic arc artifacts in the reconstruction.  The default solver
is filtered back-projection (FBP)~\cite{feldkamp1984practical} with a Ram-Lak filter, supplemented
by iterative SART for comparison.

\paragraph{CT detector offset configuration.}
Parallel-beam CT is simulated with 360 projection angles over $[0, \pi)$
and 128 detector elements.  The object is a $128 \times 128$ Shepp--Logan
phantom.  Detector offset mismatch is simulated by shifting the sinogram
along the detector axis by $\Delta s \in \{2, 5, 10, 20\}$\,pixels.  The
default solver is filtered backprojection (FBP) with Shepp--Logan filter.
Calibration uses grid search over 21 offset candidates in $[-25, 25]$\,px,
selecting the value that maximizes reconstruction PSNR.

\paragraph{Ptychography configuration.}
Electron ptychography uses 4D-STEM data: a focused probe is scanned across
a $128 \times 128$ grid, recording a convergent-beam electron diffraction
(CBED) pattern at each position (300\,kV, SrTiO$_3$ [001] from
Zenodo~5113449).  The forward model is Fresnel near-field propagation of
the exit wave through the specimen.  Mismatch is parameterized as probe
position jitter $\Delta p \in \{0.5, 1.0, 2.0, 4.0\}$\,px, simulating
scan-coil hysteresis or sample drift.  The default solver is integrated
centre-of-mass (iCoM) phase retrieval.  Calibration uses ePIE blind
position correction~\cite{maiden2009ptychography}.

\paragraph{Cryo-EM configuration.}
Cryo-EM imaging is modeled as contrast transfer function (CTF) filtering
of a 2D projected potential ($128 \times 128$, pixel size $0.1$\,nm).  The
CTF is parameterized by defocus $\Delta f = -2000$\,nm, spherical aberration
$C_s = 2.0$\,mm, electron wavelength $\lambda = 0.00251$\,nm (300\,kV),
with a B-factor envelope ($B = 2$\,nm$^2$) and ice-thickness attenuation
(mean free path $350$\,nm, thickness $50$\,nm).  Mismatch is parameterized
as defocus error $\Delta(\Delta f) \in \{100, 200, 500, 1000\}$\,nm.  The
default solver is Wiener filtering with SNR\,$= 50$.  Calibration uses grid
search over 51 defocus values in $[-3000, -500]$\,nm, selecting the value
that maximizes reconstruction PSNR.

\paragraph{MRI configuration.}
Cartesian $k$-space sampling with 4$\times$ acceleration (25\% of
$k$-space lines acquired).  Mismatch is parameterized as a 5\% multiplicative
error in the coil sensitivity maps used for parallel imaging reconstruction.
The default solver is CG-SENSE~\cite{pruessmann1999sense} with $\ell_1$-wavelet regularization ($\lambda = 10^{-3}$, 30 iterations).
\textit{Scenario~IV calibration (ESPIRiT-adopted).}
For MRI, Scenario~IV adopts ESPIRiT~\cite{uecker2014espirit} as the calibration method: coil sensitivity maps are estimated directly from the central $k$-space autocalibration signal (ACS) region via eigenvalue decomposition of the calibration matrix, then used in the CG-SENSE reconstruction.  When ACS data is sufficient ($\geq 48$ lines), ESPIRiT achieves 85--95\% of the oracle correction ceiling.  For data-limited regimes ($< 32$ ACS lines) where ESPIRiT degrades, the fallback is a grid search over per-coil amplitude and phase scaling (beam-search); this fallback is also the Sc.~IV method for all modalities without an established specialist calibration algorithm (CASSI, CACTI, compressive holography, ultrasound).  The protocol evaluates three conditions---Sc.~I (true maps), Sc.~II (5\% mismatched maps), and Sc.~IV (ESPIRiT-calibrated maps)---on the same CG-SENSE solver to isolate calibration quality from solver design.

\paragraph{Ultrasound configuration.}
Pulse-echo B-mode imaging is simulated with a 64-element linear array,
5\,MHz center frequency, element pitch $0.3$\,mm, sampling rate 100\,MHz
with 1024 time samples, and speed of sound (SoS) $c = 1540$\,m/s.  The
tissue phantom is a $128 \times 128$ reflectivity map with
Rayleigh-distributed speckle, point scatterers, and anechoic cyst regions.
Frequency-dependent attenuation follows
$\alpha(f) = 0.5\,\text{dB\,cm}^{-1}\text{MHz}^{-1}$.  Mismatch is
parameterized as SoS error $\Delta c \in \{50, 100, 200, 400\}$\,m/s,
perturbing time-of-flight calculations.  The default solver is delay-and-sum
(DAS) beamforming (operator adjoint).  Calibration uses grid search over
21 SoS candidates in $[1440, 1640]$\,m/s, selecting the value that
maximizes reconstruction PSNR.

\paragraph{Controlled hardware experiment protocol.}
The software-perturbation protocol above applies calibrated mask shifts to
existing real measurements.  A full hardware-in-the-loop validation
requires physically displacing the coded aperture mask and re-acquiring
data.  The protocol proceeds as follows: (i)~acquire a baseline dataset
with the mask at its factory-calibrated position; (ii)~physically translate
the mask by a known displacement ($\Delta x \in \{0.25, 0.5, 1.0\}$\,px
equivalent, verified by micrometer stage) and re-acquire under identical
illumination; (iii)~reconstruct both datasets with the factory mask
specification and compute the PSNR degradation and measurement residual;
(iv)~apply \pwm{} autonomous calibration and measure recovery.  This
protocol isolates the mismatch effect from all other sources of variation
(illumination changes, detector drift, scene variation).  Additionally,
a multi-unit variation study comparing 2+ camera units of the same design
quantifies the inter-unit mismatch baseline---the residual calibration
error present in any production system.  Clinical CT phantom validation
(ACR Gammex 464) and clinical MRI validation (fastMRI~\cite{zbontar2018fastmri})
configurations are detailed in Supplementary Note~S25.

\paragraph{Real experimental data.}
Three of the five multi-phantom modalities are validated on real
experimental data from established benchmarks: ultrasound uses PICMUS
experimental phantom recordings and in-vivo carotid
scans~\cite{liebgott2016picmus} plus one DeepUS CIRS-040GSE
acquisition~\cite{khan2020deepus}; cryo-EM uses five EMDB 3D density maps
(EMD-5778, EMD-2984, EMD-6287, EMD-11103, EMD-21375) projected at random
orientations; CT uses three LoDoPaB-CT patient
slices~\cite{leuschner2021lodopab}, one FIPS walnut \textmu CT
scan~\cite{hamalainen2015fips}, and one HTC\,2022
sample~\cite{meaney2022htc}.  Compressive holography and fluorescence
microscopy retain synthetic multi-phantom benchmarks as no public datasets
with calibrated ground truth are available for these modalities.

%% =========================================================================
\subsection*{Statistical Analysis}
\label{sec:methods:statistics}

\paragraph{Per-scene reporting.}
All metrics are reported per scene, not merely as dataset averages.  This
enables identification of scene-dependent failure modes (\eg{}, spectrally
flat scenes that are inherently harder for CASSI, or textureless regions
that challenge SPC).

\paragraph{Summary statistics.}
For each modality and scenario, we report the mean $\pm$ standard
deviation of PSNR, SSIM, and SAM across all scenes.  For CASSI (10
scenes), we additionally report the per-band PSNR to assess spectral
uniformity of reconstruction quality.

\paragraph{Recovery ratio confidence intervals.}
The recovery ratio $\recoveryratio$ (\cref{eq:recovery_ratio}) is a ratio
of differences and therefore sensitive to noise in the constituent PSNR
values.  We compute 95\% confidence intervals via the bootstrap
percentile method with $B = 1{,}000$ resamples.  At each bootstrap
iteration, we resample the scene set with replacement, recompute the
mean PSNR for each scenario, and derive $\recoveryratio$.  The 2.5th
and 97.5th percentiles of the bootstrap distribution define the 95\% CI.

\paragraph{Parameter recovery accuracy.}
For mismatch correction experiments, we report the root-mean-square error
(RMSE) between the estimated and true mismatch parameters:
\begin{equation}
  \text{RMSE}_k = \sqrt{\frac{1}{N_{\text{scene}}}
    \sum_{i=1}^{N_{\text{scene}}}
    (\hat\theta_{k,i} - \theta_{k,\text{true}})^2}
\end{equation}
where $k$ indexes the mismatch parameter, $i$ indexes the scene, and
$N_{\text{scene}}$ is the number of test scenes.  Uncertainty in the RMSE
is estimated via bootstrap ($B = 1{,}000$).

\paragraph{Ablation significance.}
Ablation studies (removal of the photon-budget, recoverability, or
mismatch scoring module, or RunBundle discipline) are evaluated by comparing the
full-pipeline PSNR against each ablated variant.  We report the PSNR
difference $\Delta\psnr$ per modality and verify that each component
contributes $\geq 0.5$\,dB across all validated modalities, establishing
practical significance.


%% =========================================================================
\subsection*{Code and Data Availability}
\label{sec:methods:availability}

\paragraph{Source code.}
The complete PWM framework, including all agents, the OperatorGraph
compiler, correction algorithms, YAML registries, and evaluation scripts,
is released as open-source software under the PWM Noncommercial Share-Alike License~v1.0 at
\url{https://github.com/integritynoble/Physics_World_Model}.  The
codebase is organized into two Python packages:
\texttt{pwm\_core} (core framework, agents, graph compiler, calibration
algorithms) and \texttt{pwm\_AI\_Scientist} (automated experiment
generation and analysis).

\paragraph{Reconstruction data.}
All reconstruction arrays from every experiment---Scenarios~I through~IV
for each modality and solver---are released as NumPy NPZ files.
Files are stored using Git LFS and require \texttt{allow\_pickle=True}
for loading.  Data ranges are standardized: CASSI and CACTI
reconstructions are normalized to $[0, 1]$; SPC reconstructions are in
$[0, 255]$.

\paragraph{Experiment manifests.}
Every experiment is recorded in a RunBundle v0.3.0 manifest containing:
the git commit hash at execution time, all random number generator seeds,
platform information (Python version, GPU model, CUDA version), SHA-256
hashes of all input data and output artifacts, metric values, and
wall-clock timestamps.  These manifests enable exact reproduction of every
reported result.

\paragraph{Registry data.}
All 9 YAML registries (7,034 lines total) that drive the agent system---including
modality definitions, graph templates, photon models, mismatch databases,
compression tables, solver configurations, primitive specifications,
dataset paths, and acceptance thresholds---are publicly available in the
repository under \texttt{packages/pwm\_core/contrib/}.
The \texttt{ExperimentSpec} JSON schemas used for pipeline input
validation are included alongside worked examples in
\texttt{examples/}.
 from the main manuscript.

\section*{Online Methods}

%% =========================================================================
\subsection*{OperatorGraph Specification}
\label{sec:methods:operatorgraph}

\paragraph{Formal definition.}
The \opg{} intermediate representation encodes the forward physics of any
computational imaging modality as a directed acyclic graph (DAG)
$\mathcal{G} = (\mathcal{V}, \mathcal{E})$.
Each node $v_i \in \mathcal{V}$ wraps a \emph{primitive operator} and
implements two entry points: $\texttt{forward}(x) \to y$ and
$\texttt{adjoint}(y) \to x$, the latter defined only when the primitive
is linear.  Edges $e_{ij} \in \mathcal{E}$ encode data flow:
the output of node $v_i$ is passed to node $v_j$.
Each node additionally exposes a set of learnable parameters
$\boldsymbol{\theta}_i$ that may be perturbed during mismatch simulation
or optimized during calibration, as well as read-only metadata flags
(\texttt{is\_linear}, \texttt{is\_stochastic}, \texttt{is\_differentiable}).
The graph is stored as a declarative YAML specification
(\texttt{OperatorGraphSpec}) and compiled to an executable
\texttt{GraphOperator} object by the \texttt{GraphCompiler}.

\paragraph{Node types.}
Primitive operators fall into two categories:

\begin{itemize}
  \item \textbf{Linear operators.}  Convolution (\texttt{conv2d}), mask
    modulation (\texttt{mask\_modulate}), sub-pixel shift (\texttt{subpixel\_shift\_2d}),
    Radon transform (\texttt{radon\_fanbeam}),
    Fourier encoding (\texttt{fourier\_encode}), spectral dispersion
    (\texttt{spectral\_disperse}), Fresnel propagation
    (\texttt{fresnel\_propagate}), random projection
    (\texttt{random\_project}), and structured illumination
    (\texttt{sim\_modulate}).
    Each implements both \texttt{forward()} and \texttt{adjoint()}.
  \item \textbf{Nonlinear operators.}  Squared magnitude
    (\texttt{magnitude\_sq}), Poisson--Gaussian noise
    (\texttt{poisson\_gaussian}), saturation clipping
    (\texttt{saturation\_clip}), phase retrieval nonlinearity
    (\texttt{phase\_abs}), and detector quantization
    (\texttt{quantize}).
    These set \texttt{is\_linear = False} and raise
    \texttt{NotImplementedError} on \texttt{adjoint()}, except where a
    well-defined pseudo-adjoint exists (\eg{}, the identity adjoint for
    magnitude-squared in Gerchberg--Saxton-type algorithms).
\end{itemize}

\paragraph{Adjoint consistency and Lipschitz bounds as mathematical safety rails.}
Two structural constraints distinguish the \opg{} IR from standard deep-learning
forward models and prevent either nonlinear primitive family from becoming a
universal approximator.

First, \emph{adjoint consistency}: correctness of every linear primitive
is verified by a randomized dot-product test.  For a primitive $A$ with
forward map $A:\mathbb{R}^n \to \mathbb{R}^m$, we draw $x \sim \mathcal{N}(0,I_n)$
and $y \sim \mathcal{N}(0,I_m)$ and compute
\begin{equation}
  \delta = \frac{|\langle A^* y,\, x \rangle - \langle y,\, Ax \rangle|}
               {\max(|\langle A^* y,\, x \rangle|,\; \epsilon)}
  \label{eq:adjoint_check}
\end{equation}
where $\epsilon = 10^{-12}$ guards against division by zero.  The test
is repeated $n_{\text{trials}} = 5$ times with independent random draws;
the primitive passes if $\delta_{\max} < 10^{-6}$.  At the graph level,
a compiled \texttt{GraphOperator} composed entirely of linear nodes
executes the same test over the composed forward--adjoint chain.  A
\texttt{GraphAdjointCheckReport} records $n_{\text{trials}}$,
$\delta_{\max}$, and $\bar\delta$ for audit.  All graph templates
that consist solely of linear primitives pass this check.
Adjoint consistency is not merely a numerical correctness guarantee: it
is a \emph{physical faithfulness certificate}.  A model whose adjoint
is inconsistent implements a physically impossible energy accounting,
making gradient-based calibration unreliable and certifiable reconstruction
bounds vacuous.  Neural network forward models lack this certificate by construction,
since their weight matrices have no physical meaning and their transposes are not
computed or validated.

Second, \emph{Lipschitz bounds}: the two nonlinear primitive families,
Detect~$D$ and Transform~$\Lambda$, are each restricted to five canonical
function families with at most two scalar parameters.  This restriction
enforces bounded Lipschitz constants ($\mathrm{Lip}(\Lambda_k) \leq L$ for
the explicit $L$ derived in~\cite{yang2026fpt}) and prevents either primitive
from becoming a universal approximator in the sense of Cybenko~\cite{cybenko1989approximation}.
Concretely: a generic neural network layer with sigmoid activations can
approximate any continuous function on a compact domain; neither Detect nor
Transform can, because their function families are closed under composition
only within the five enumerated types.  This non-universality is the key
property that makes the 11-primitive library \emph{falsifiable}: if a
modality's forward physics cannot be represented within the restricted
families, the extension protocol (Main Text, ``Extension protocol'') fires
and a new primitive must be justified physically.  The combination of
adjoint consistency and bounded Lipschitz nonlinearity provides the
mathematical safety rails that make \opg{}-based reconstruction both
interpretable and certifiable---properties that standard physics-informed
neural networks~\cite{raissi2019pinn} achieve only approximately and
without modality-agnostic guarantees.

\paragraph{Graph compilation.}
The compiler executes a four-stage pipeline:
\begin{enumerate}
  \item \textbf{Validate.}  Confirm acyclicity via topological sort
    (Kahn's algorithm), verify that every \texttt{primitive\_id}
    exists in the global \texttt{PRIMITIVE\_REGISTRY}, reject
    duplicate \texttt{node\_id} values, and optionally verify shape
    compatibility along edges when a \texttt{canonical\_chain}
    metadata flag is set.
  \item \textbf{Bind.}  Instantiate each primitive with its parameter
    dictionary $\boldsymbol{\theta}_i$.
  \item \textbf{Plan forward.}  The topological sort yields a sequential
    execution plan $(v_{\pi(1)}, \ldots, v_{\pi(|\mathcal{V}|)})$.
  \item \textbf{Plan adjoint.}  For graphs where
    \texttt{all\_linear = True}, the adjoint plan reverses the
    topological order and applies each node's individual adjoint in
    sequence, implementing the chain rule $A^* = A_1^*
    \circ \cdots \circ A_{|\mathcal{V}|}^*$ for a composition $A = A_{|\mathcal{V}|}
    \circ \cdots \circ A_1$.  For graphs containing nonlinear nodes,
    the adjoint plan is not generated, and any call to
    \texttt{adjoint()} raises \texttt{NotImplementedError} at runtime.
\end{enumerate}
The compiled \texttt{GraphOperator} is serializable to JSON and hashable
via SHA-256 for provenance tracking in RunBundle manifests.

\paragraph{Template library.}
The \texttt{graph\_templates.yaml} registry contains 170 templates
spanning all five carrier families. Of these, 12 modalities carry full
end-to-end Scenario~I--IV correction validation across all five carrier
families (Supplementary Table~S3). Representative modalities by carrier family include:

\begin{itemize}
  \item \textbf{Optical and X-ray photons:}  CASSI, SPC, CACTI, structured
    illumination microscopy (SIM), confocal, light-sheet, holography,
    ptychography, Fourier ptychographic microscopy (FPM), optical coherence
    tomography (OCT), lensless imaging, light field, integral imaging,
    neural radiance fields (NeRF), Gaussian splatting, fluorescence
    lifetime imaging (FLIM), diffuse optical tomography (DOT),
    phase retrieval, X-ray computed tomography (CT), and cone-beam CT
    (CBCT).
  \item \textbf{Electrons:}  Electron diffraction, electron backscatter
    diffraction (EBSD), electron energy loss spectroscopy (EELS), and
    electron holography.
  \item \textbf{Nuclear spins (MRI):}  Functional MRI (fMRI), diffusion-weighted
    MRI (DW-MRI), and magnetic resonance spectroscopy (MRS).
  \item \textbf{Acoustic waves:}  Ultrasound B-mode, Doppler ultrasound,
    shear-wave elastography, sonar, and photoacoustic tomography
    (combines optical excitation with acoustic detection).
  \item \textbf{Particle-beam probes:}  Neutron tomography, proton CT,
    and muon tomography.  Each decomposes into existing primitives with
    no extension required: neutron CT follows the X-ray CT DAG
    ($P \!\to\! \Lambda \!\to\! D$, Beer--Lambert attenuation with nuclear
    cross-section $\mu_n$); proton CT inserts a Bethe--Bloch
    stopping-power stage as Transform~$\Lambda$ (\#11); muon tomography
    uses Scatter~$R$ (\#10) for multiple Coulomb scattering and
    POCA vertex reconstruction (Supplementary Note~S24).
\end{itemize}

\paragraph{Fidelity levels.}
Each template is parameterized by a fidelity level that controls the
degree of physical realism in the simulated forward model:

\begin{description}
  \item[Level~1 (Linear, shift-invariant):]  The forward model is a
    linear, spatially uniform operator---the simplest approximation,
    suitable for initial diagnostics and rapid prototyping.
  \item[Level~2 (Linear, shift-variant):]  Spatially varying operator
    parameters (e.g.\ non-uniform illumination, position-dependent PSF,
    multi-coil sensitivity maps in MRI).  Adds a modality-appropriate
    noise model (Poisson shot noise plus Gaussian read noise for
    photon-counting modalities, Rician noise for MRI, Poisson for CT).
  \item[Level~3 (Nonlinear, ray/wave-based):]  Includes nonlinear
    effects such as beam hardening, phase wrapping, and scattering,
    captured by the Transform~$\Lambda$ primitive.
    Perturbation families and ranges are specified in
    \texttt{mismatch\_db.yaml}.
  \item[Level~4 (Full-wave / Monte Carlo):]  Complete physical
    simulation including wave-optical propagation, spatially varying
    aberrations, detector nonlinearities, and environmental drift.
    Currently implemented for holography and ptychography.
\end{description}

\noindent All 11 primitives in the library $\Blib$ apply uniformly across every fidelity level; the levels organize simulation complexity, not theorem scope.


%% =========================================================================
\subsection*{Triad Decomposition Formalization}
\label{sec:methods:triad}

The \triad{} asserts that the quality of any computational imaging
reconstruction is bounded by three fundamental gates.  Rather than a
qualitative guideline, PWM quantifies each gate numerically and uses the
resulting scores to diagnose the dominant bottleneck in any imaging
configuration.

\paragraph{Gate~1 (Recoverability).}
Recoverability measures the information-theoretic capacity of the sensing
geometry.  We quantify it via the \emph{effective compression ratio}
$r = m / n$, where $m$ is the number of independent measurements and $n$
the dimension of the signal.  The \texttt{compression\_db.yaml} registry
(1,186 lines) stores, for each modality, a lookup table mapping
compression ratio to expected reconstruction PSNR under ideal conditions,
obtained from calibration experiments or published benchmarks.  Each entry
carries a \texttt{provenance} field citing the source (paper DOI, internal
experiment ID, or theoretical formula).  Additional recoverability
indicators include the effective rank of the measurement matrix (estimated
via randomized SVD for large operators), the dimension of the null space,
and the restricted isometry property (RIP) constant where analytically
tractable (\eg{}, for Gaussian random projections in SPC).

\paragraph{Gate~2 (Carrier Budget).}
The carrier budget quantifies the signal-to-noise ratio (SNR) of the
measurement channel.  The photon-budget module consumes the
\texttt{photon\_db.yaml} registry which stores, per modality,
a deterministic photon model parameterized by source power, quantum
efficiency, exposure time, and detector characteristics.  It
classifies the noise regime into one of three categories:
\emph{shot-limited} (Poisson-dominated, SNR $\propto \sqrt{N_{\text{photon}}}$),
\emph{read-limited} (Gaussian read noise dominates, SNR $\propto
N_{\text{photon}} / \sigma_{\text{read}}$), and \emph{dark-current-limited}
(long exposures where dark current accumulation dominates).  The output is
a \texttt{PhotonReport} containing the estimated SNR in decibels, the noise
regime classification, per-element photon count, and a feasibility verdict
(\texttt{sufficient}, \texttt{marginal}, or \texttt{insufficient}).

\paragraph{Gate~3 (Operator Mismatch).}
Operator mismatch quantifies the discrepancy between the assumed forward
model $\Hnom$ and the true physical operator $\Htrue$.  The
mismatch scoring module consults \texttt{mismatch\_db.yaml}
which catalogs, for each modality, the set of mismatch parameters
(spatial shifts, rotational offsets, dispersion errors, PSF deviations,
coil sensitivity errors, center-of-rotation offsets, \etc{}), their
typical ranges, and available correction methods.  The mismatch severity
score $s \in [0, 1]$ is computed as the normalized $\ell_2$ distance
$\| \boldsymbol{\theta}_{\text{true}} - \boldsymbol{\theta}_{\text{nom}} \|
/ \| \boldsymbol{\theta}_{\text{range}} \|$, where
$\boldsymbol{\theta}_{\text{range}}$ is the per-parameter dynamic range
from the registry.  Sensitivity analysis $\partial \psnr / \partial
\theta_k$ is estimated via finite differences on the forward model.  The
output is a \texttt{MismatchReport} containing the severity score, the
dominant mismatch parameter, the recommended correction method, and the
expected PSNR gain from correction.

\paragraph{Gate binding determination.}
Given reconstruction results under the four-scenario protocol
(the Evaluation Protocol section below), PWM identifies the dominant gate by
comparing three cost terms:
\begin{align}
  C_{\text{mismatch}} &= \psnr_{\text{I}} - \psnr_{\text{II}} \label{eq:cost_mismatch} \\
  C_{\text{noise}}    &= \psnr_{\text{ideal}} - \psnr_{\text{noisy}} \label{eq:cost_noise} \\
  C_{\text{recover}}  &= \psnr_{\text{limit}} - \psnr_{\text{I}} \label{eq:cost_recover}
\end{align}
where $\psnr_{\text{I}}$ is the reconstruction PSNR under Scenario~I
(ideal operator), $\psnr_{\text{II}}$ under Scenario~II (mismatched
operator), $\psnr_{\text{noisy}}$ under the corresponding noisy condition,
and $\psnr_{\text{limit}}$ is the theoretical upper bound from the
compression table.  The dominant gate is $\arg\max_g C_g$.

\paragraph{TriadReport.}
For every diagnosis, \pwm{} produces a \triadreport{}: a Pydantic-validated
structured artifact comprising
\texttt{dominant\_gate} (enum: \texttt{recoverability},
\texttt{carrier\_budget}, \texttt{operator\_mismatch}),
\texttt{evidence\_scores} (three floats, one per gate),
\texttt{confidence\_interval} (float, 95\% CI width from bootstrap),
\texttt{recommended\_action} (string, \eg{} ``increase compression ratio''
or ``apply mismatch correction''), and
\texttt{parameter\_sensitivities} (dictionary mapping each mismatch
parameter name to its $\partial \psnr / \partial \theta_k$ value).
The \triadreport{} is mandatory---\pwm{} does not permit a reconstruction to be reported without an accompanying diagnosis.

\paragraph{Recovery ratio.}
We define the \emph{recovery ratio}
\begin{equation}
  \recoveryratio = \frac{\psnr_{\text{IV}} - \psnr_{\text{II}}}
                        {\psnr_{\text{I}} - \psnr_{\text{II}}}
  \label{eq:recovery_ratio}
\end{equation}
which lies in $[0, 1]$ under standard convexity conditions (see Supplementary Note~1 for formal analysis; values $\recoveryratio > 1$ are possible when the corrected operator provides beneficial regularization).
$\recoveryratio = 0$ indicates that calibration yields no benefit (mismatch
is not the bottleneck), while $\recoveryratio = 1$ indicates that
calibration fully closes the mismatch gap.


%% =========================================================================
\subsection*{Agent System Architecture}
\label{sec:methods:agents}

The three gates are implemented as a deterministic scoring pipeline:
an orchestrator maps the user request to a modality from the registry,
then dispatches three parallel scoring modules---one per gate---whose
reports are fused by a bottleneck classifier that identifies the dominant
gate and by a negotiator that enforces cross-gate veto rules
(e.g., rejecting reconstruction when photon starvation coincides with
aggressive compression).  All quantitative decisions are deterministic;
an LLM is used only as an optional fallback for natural-language modality
mapping and is never on the scoring path.  The full agent architecture
and ablation studies are described in a companion paper.

\paragraph{Contract system.}
Inter-agent communication uses 25 Pydantic v2 contract models.  All
contracts inherit from \texttt{StrictBaseModel}, which enforces
\texttt{extra="forbid"} (no unexpected fields),
\texttt{validate\_assignment=True} (mutations re-validated), and a model
validator that rejects NaN and Inf in any float field.  Bounded scores use
\texttt{Field(ge=0.0, le=1.0)}.  Enums are string enums for human-readable
JSON serialization.  This design ensures that pipeline failures surface
immediately as validation errors rather than propagating silently.

\paragraph{YAML registries.}
The system is driven by 9 YAML registries totalling 7,034 lines:
\texttt{modalities.yaml} (modality definitions),
\texttt{graph\_templates.yaml} (OperatorGraph skeletons),
\texttt{photon\_db.yaml} (photon models),
\texttt{mismatch\_db.yaml} (mismatch parameters and correction methods),
\texttt{compression\_db.yaml} (recoverability tables with provenance),
\texttt{solver\_registry.yaml} (solver configurations),
\texttt{primitives.yaml} (primitive operator metadata),
\texttt{dataset\_registry.yaml} (dataset locations and formats), and
\texttt{acceptance\_thresholds.yaml} (pass/fail thresholds per metric).


%% =========================================================================
\subsection*{Correction Algorithms}
\label{sec:methods:correction}

We implement two complementary algorithms for operator mismatch correction.
Crucially, both algorithms operate on the forward operator parameters
$\boldsymbol{\theta}$ rather than the reconstruction solver weights,
making them \emph{solver-agnostic}: the corrected operator
$H(\hat{\boldsymbol{\theta}})$ benefits any downstream solver (GAP-TV,
MST-L, HDNet~\cite{hu2022hdnet}, CST, \etc{}) without retraining.

\paragraph{Algorithm~1: Hierarchical Beam Search.}
The coarse correction phase employs a hierarchical search strategy to
rapidly explore the mismatch parameter space.  For CASSI, the five-parameter
mismatch model comprises mask affine parameters (spatial shifts $dx$, $dy$
and rotation $\theta$) and dispersion parameters (slope $a_1$ and axis
angle $\alpha$); an optional sixth parameter, PSF width $\sigma_{\text{psf}}$, is available but not used in the primary experiments.
The algorithm proceeds as follows:

\begin{enumerate}
  \item \textbf{1D sweeps.}  Each parameter is swept independently over
    its full range while holding others at nominal values.  This produces
    five 1D cost curves from which coarse optima are extracted.
  \item \textbf{3D beam search.}  The mask affine subspace
    $(dx, dy, \theta)$ is searched over a $5 \times 5 \times 5$ grid
    centered on the 1D optima.  The top-$k$ ($k = 5$) candidates by
    reconstruction PSNR are retained.
  \item \textbf{2D beam search.}  For each retained mask candidate,
    the dispersion subspace $(a_1, \alpha)$ is searched over a
    $5 \times 7$ grid.  The joint top-$k$ candidates are retained.
  \item \textbf{Coordinate descent refinement.}  Three rounds of
    univariate refinement on each parameter, shrinking the search
    interval by factor~2 at each round, produce the final estimate
    $\hat{\boldsymbol{\theta}}_{\text{Alg1}}$.
\end{enumerate}

Total runtime is approximately 300 seconds per scene on a single GPU.
Accuracy is $\pm 0.1$--$0.2$ pixels for spatial parameters and
$\pm 0.05^\circ$ for angular parameters.

\paragraph{Algorithm~2: Joint Gradient Refinement.}
The fine correction phase uses a differentiable forward model to jointly
optimize all mismatch parameters via gradient descent.  The key components
are:

\begin{enumerate}
  \item \textbf{Differentiable mask warp.}  The binary mask is warped by a
    continuous affine transformation using bilinear interpolation,
    implemented as a custom PyTorch module
    (\texttt{DifferentiableMaskWarpFixed}).  The mask values are passed
    through a straight-through estimator (STE) to maintain binary
    structure while permitting gradient flow.
  \item \textbf{Differentiable forward model.}  The CASSI forward model
    $y = \texttt{CASSI}(x;\, \boldsymbol{\theta})$ is implemented as a
    differentiable PyTorch module (\texttt{DifferentiableCassiForwardSTE})
    that accepts mismatch parameters as differentiable inputs.
  \item \textbf{GPU grid initialization.}  A full-range 3D grid search
    over $(dx, dy, \theta)$ with $9 \times 9 \times 7 = 567$ points
    provides diverse starting candidates.  The top 9 candidates seed
    multi-start gradient refinement.
  \item \textbf{Staged gradient refinement.}  Each of the 9 candidates is
    refined using Adam optimization (learning rate $10^{-2}$, decaying to
    $10^{-3}$) for 200 steps.  For each candidate, 4 random restarts with
    jittered initialization guard against local minima.  The loss function
    is the negative PSNR computed via an unrolled $K$-iteration
    differentiable GAP-TV solver (\texttt{DifferentiableGAPTV}, $K=10$
    unrolled iterations).
\end{enumerate}

Total runtime for Algorithm~2 is approximately 3,200 seconds
(200~steps $\times$ 4~restarts $\times$ 9~candidates with early stopping).
Accuracy improves to $\pm 0.05$--$0.1$ pixels, a 3--5$\times$ improvement
over Algorithm~1.  The two algorithms are used sequentially in practice:
Algorithm~1 provides a warm start, and Algorithm~2 refines to sub-pixel
precision.


%% =========================================================================
\subsection*{Evaluation Protocol}
\label{sec:methods:eval}

\paragraph{Four-Scenario Protocol.}
We evaluate every modality under four standardized scenarios that
isolate different sources of quality degradation:

\begin{description}
  \item[Scenario~I (Ideal):]
    $\yobs = \Htrue\, \xgt$; reconstruct with $\Htrue$.
    In this scenario the system is perfectly calibrated ($\Htrue = \Hnom$),
    so the operator used for reconstruction matches the one that generated
    the data.  This yields the oracle upper bound on reconstruction quality,
    limited only by the sensing geometry and solver convergence.
  \item[Scenario~II (Mismatch):]
    $\yobs = \Htrue\, \xgt$; reconstruct with $\Hnom$
    ($\Hnom \neq \Htrue$).
    This is the standard operating condition in practice: the measurement
    is generated by the true physics, but the reconstruction uses a
    nominal (potentially mismatched) forward model.
  \item[Scenario~III (Oracle):]
    Same measurements as Scenario~II
    ($\yobs = \Htrue\, \xgt$ with $\Htrue \neq \Hnom$);
    reconstruct with $\Htrue$ instead of $\Hnom$.
    Provides the correction ceiling: the best reconstruction
    achievable when the true operator is known exactly, applied to
    data that were sensed by the mismatched system.  The gap between
    Scenario~III and Scenario~I reveals the irreducible loss from
    the degraded sensing configuration itself (e.g., a shifted mask
    pattern is suboptimal even when perfectly known).
  \item[Scenario~IV (Corrected):]
    $\yobs = \Htrue\, \xgt$; reconstruct with
    $\hat{H} = H(\hat{\boldsymbol{\theta}})$ where
    $\hat{\boldsymbol{\theta}}$ is estimated by Algorithms~1 and~2.
    This quantifies the benefit of mismatch calibration.
\end{description}

\paragraph{Metrics.}
Reconstruction quality is assessed using three complementary metrics:
\begin{itemize}
  \item \textbf{PSNR} (peak signal-to-noise ratio, in dB): the primary
    metric, computed per scene and averaged.  For signals normalized to
    $[0, 1]$, $\psnr = 10 \log_{10}(1 / \text{MSE})$.  For SPC data
    normalized to $[0, 255]$, the peak value is 255.
  \item \textbf{SSIM} (structural similarity index): captures perceptual
    quality including luminance, contrast, and structural components,
    computed with a Gaussian window of width 11 and standard deviation
    1.5.
  \item \textbf{SAM} (spectral angle mapper): for hyperspectral
    modalities (CASSI), measures the angle between predicted and true
    spectral vectors at each spatial location, reported in degrees.
    Lower is better.
\end{itemize}

\paragraph{Datasets.}
\begin{itemize}
  \item \textbf{CASSI:} 10 scenes from the KAIST dataset~\cite{choi2017kaist},
    each a $256 \times 256 \times 28$ spectral cube (28 spectral bands
    from 450\,nm to 650\,nm).  Data range $[0, 1]$.
  \item \textbf{CACTI:} 6 benchmark videos, each $256 \times 256 \times 8$
    (8 temporal frames encoded per snapshot).  Data range $[0, 1]$.
  \item \textbf{SPC:} 11 natural images from the Set11 benchmark, each
    $256 \times 256$ grayscale.  Data range $[0, 255]$.
\end{itemize}

\noindent All per-scene metrics are reported individually as well as
averaged, and all reconstruction arrays are saved as NumPy NPZ files.


%% =========================================================================
\subsection*{Experimental Details}
\label{sec:methods:experimental}

\paragraph{Hardware.}
All experiments are conducted on a single NVIDIA GPU.  Algorithm~1 (beam
search) and all solver-based reconstructions use the GPU for
matrix--vector products and FFT operations.  Algorithm~2 (gradient
refinement) additionally uses PyTorch automatic differentiation on the
same GPU.

\paragraph{CASSI configuration.}
The coded aperture snapshot spectral imaging (CASSI) system uses a TSA-Net
binary mask of dimensions $256 \times 256$, with 28 spectral bands
dispersed along the spatial dimension.  The five-parameter mismatch model
$\mismatchparam = (dx, dy, \theta, a_1, \alpha)$
describes: mask spatial shift in $x$ ($dx$, pixels), mask spatial shift in
$y$ ($dy$, pixels), mask rotation angle ($\theta$, degrees), dispersion
slope ($a_1$, pixels per band), and dispersion axis angle ($\alpha$,
degrees).  An optional sixth parameter, PSF blur width ($\sigma_{\text{psf}}$, pixels), is available but not used in the primary experiments.
These mismatch parameter values were determined through systematic characterization of realistic CASSI assembly errors (Supplementary Note~8). The true mismatch parameters are
$\mismatchparam_{\text{true}} = (dx = 0.5\;\text{px},\; dy = 0.3\;\text{px},\;
\theta = 0.1^\circ,\; a_1 = 2.02,\;
\alpha = 0.15^\circ)$.
Solvers evaluated include TwIST~\cite{bioucasdias2007twist}, GAP-TV~\cite{yuan2016gaptv}, DGSMP~\cite{huang2021dgsmp}, MST-L~\cite{cai2022mst}, and CST-L~\cite{cai2022cst}, all of which receive the same operator and differ only in their reconstruction algorithm. The supplementary per-scene analysis additionally includes DeSCI~\cite{liu2019rank} and HDNet~\cite{hu2022hdnet}.
\textit{Real-data validation.}
The TSA real hyperspectral dataset~\cite{wagadarikar2008cassi} consists of
5 scenes at $660 \times 660$ spatial resolution with 28 spectral bands
and mask-shift step~2.  Four solvers are evaluated: GAP-TV (200 iterations),
HDNet (pre-trained checkpoint, full spatial resolution), MST-S and MST-L
(pre-trained checkpoints, centre-cropped to $256 \times 256$ due to
hardcoded spatial assumptions in the model architecture).  The coded
aperture mask is perturbed by $dx = 0.5$\,px, $dy = 0.3$\,px to simulate
assembly-induced mismatch.  No ground truth is available; quality is
assessed via the normalised measurement residual
$r = \|\mathbf{y} - H\hat{\mathbf{x}}\|^2 / \|\mathbf{y}\|^2$.

\paragraph{CACTI configuration.}
The coded aperture compressive temporal imaging system uses binary
temporal masks of dimensions $256 \times 256$, encoding 8 video frames
into a single snapshot measurement.  Mismatch is parameterized as a
temporal mask timing offset (sub-frame shift).  The default solver is
EfficientSCI~\cite{wang2023efficientsci}.
\textit{Real-data validation.}
The EfficientSCI real temporal dataset~\cite{wang2023efficientsci} consists
of 4 dynamic scenes (duomino, hand, pendulumBall, waterBalloon) at
$512 \times 512$ with compression ratio~10.  The real mask is stored
separately from the measurement data.  Two solvers are evaluated:
GAP-TV (50 iterations) and PnP-FFDNet (50 iterations with FFDNet denoiser).
Mismatch is induced by shifting the mask by $dx = 0.5$\,px, $dy = 0.3$\,px.
Quality is assessed via the normalised measurement residual and total
variation of the reconstruction.

\paragraph{SPC configuration.}
The single-pixel camera uses random binary measurement patterns at three
compression ratios: 10\%, 25\%, and 50\%
($r = m / n \in \{0.10, 0.25, 0.50\}$).  Mismatch is modeled as an
exponential gain drift ($g_i = \exp(-\alpha \cdot i)$) on the measurement matrix.  The default solver
is FISTA-TV with total-variation regularization.

\paragraph{Lensless imaging configuration.}
Lensless cameras replace conventional optics with a diffuser placed directly
on the sensor, making the forward model a 2D convolution with the measured
point-spread function (PSF).  The object is a $64 \times 64$ image
convolved with a Gaussian PSF of width $\sigma = 2.0$\,px.  Mismatch is
parameterized as PSF width error
$\Delta\sigma \in \{0.5, 1.0, 2.0, 3.0\}$\,px, simulating an error in
the assumed pinhole-to-sensor distance.  Five structurally diverse phantoms
are used ($N{=}5$).  The default solver is FISTA-TV (80 iterations).
Calibration uses forward-model fitting on a known reference target:
$\sigma_{\mathrm{cal}} = \mathrm{argmin}_\sigma \|y_{\mathrm{cal}} -
h_\sigma * x_{\mathrm{cal}}\|^2$.
\textit{Real-data validation.}
Five DiffuserCam DLMD Mirflickr scenes~\cite{Antipa2018} are used to
corroborate simulation results on physical photographic content.

\paragraph{Compressive holography configuration.}
Multi-depth inline holography encodes a $(4 \times 64 \times 64)$ 3D
object into a single $(64 \times 64)$ hologram via Fresnel propagation at
four depth planes with spacing $200$\,\textmu m, wavelength
$\lambda = 532$\,nm, and pixel pitch $5$\,\textmu m.  Mismatch is
parameterized as propagation distance error
$\Delta z \in \{10, 50, 100, 200\}$\,\textmu m applied uniformly to all
depth planes.  The default solver is FISTA-TV (80 iterations,
$\lambda_{\mathrm{TV}} = 0.005$).  Calibration uses hologram residual
minimisation: grid search over distance error candidates with the search
range clamped to ensure all propagation distances remain positive.

\paragraph{Fluorescence microscopy configuration.}
Widefield fluorescence imaging uses a dual-PSF Stokes-shift model with
Gaussian excitation PSF ($\sigma_{\mathrm{ex}} = 1.5$\,px) and emission
PSF ($\sigma_{\mathrm{em}} = 2.0$\,px), quantum yield $\eta = 0.7$, and
background level $b = 0.02$.  The specimen is a $64 \times 64$ image with
fluorescent puncta, diffuse organelles, and filamentary structures.
Mismatch is parameterized as PSF sigma error
$\Delta\sigma \in \{0.1, 0.3, 0.5, 1.0\}$\,px applied to both excitation
and emission PSFs.  The default solver is Richardson--Lucy (80 iterations).
Calibration uses a $9 \times 9$ grid search over
$(\sigma_{\mathrm{ex}}, \sigma_{\mathrm{em}})$ centered at the nominal
values, selecting the pair that minimizes the measurement residual.

\paragraph{CT configuration.}
Fan-beam geometry with 180 projections over $180^\circ$.  Mismatch is
modeled as a center-of-rotation (CoR) offset, which produces
characteristic arc artifacts in the reconstruction.  The default solver
is filtered back-projection (FBP)~\cite{feldkamp1984practical} with a Ram-Lak filter, supplemented
by iterative SART for comparison.

\paragraph{CT detector offset configuration.}
Parallel-beam CT is simulated with 360 projection angles over $[0, \pi)$
and 128 detector elements.  The object is a $128 \times 128$ Shepp--Logan
phantom.  Detector offset mismatch is simulated by shifting the sinogram
along the detector axis by $\Delta s \in \{2, 5, 10, 20\}$\,pixels.  The
default solver is filtered backprojection (FBP) with Shepp--Logan filter.
Calibration uses grid search over 21 offset candidates in $[-25, 25]$\,px,
selecting the value that maximizes reconstruction PSNR.

\paragraph{Ptychography configuration.}
Electron ptychography uses 4D-STEM data: a focused probe is scanned across
a $128 \times 128$ grid, recording a convergent-beam electron diffraction
(CBED) pattern at each position (300\,kV, SrTiO$_3$ [001] from
Zenodo~5113449).  The forward model is Fresnel near-field propagation of
the exit wave through the specimen.  Mismatch is parameterized as probe
position jitter $\Delta p \in \{0.5, 1.0, 2.0, 4.0\}$\,px, simulating
scan-coil hysteresis or sample drift.  The default solver is integrated
centre-of-mass (iCoM) phase retrieval.  Calibration uses ePIE blind
position correction~\cite{maiden2009ptychography}.

\paragraph{Cryo-EM configuration.}
Cryo-EM imaging is modeled as contrast transfer function (CTF) filtering
of a 2D projected potential ($128 \times 128$, pixel size $0.1$\,nm).  The
CTF is parameterized by defocus $\Delta f = -2000$\,nm, spherical aberration
$C_s = 2.0$\,mm, electron wavelength $\lambda = 0.00251$\,nm (300\,kV),
with a B-factor envelope ($B = 2$\,nm$^2$) and ice-thickness attenuation
(mean free path $350$\,nm, thickness $50$\,nm).  Mismatch is parameterized
as defocus error $\Delta(\Delta f) \in \{100, 200, 500, 1000\}$\,nm.  The
default solver is Wiener filtering with SNR\,$= 50$.  Calibration uses grid
search over 51 defocus values in $[-3000, -500]$\,nm, selecting the value
that maximizes reconstruction PSNR.

\paragraph{MRI configuration.}
Cartesian $k$-space sampling with 4$\times$ acceleration (25\% of
$k$-space lines acquired).  Mismatch is parameterized as a 5\% multiplicative
error in the coil sensitivity maps used for parallel imaging reconstruction.
The default solver is CG-SENSE~\cite{pruessmann1999sense} with $\ell_1$-wavelet regularization ($\lambda = 10^{-3}$, 30 iterations).
\textit{Scenario~IV calibration (ESPIRiT-adopted).}
For MRI, Scenario~IV adopts ESPIRiT~\cite{uecker2014espirit} as the calibration method: coil sensitivity maps are estimated directly from the central $k$-space autocalibration signal (ACS) region via eigenvalue decomposition of the calibration matrix, then used in the CG-SENSE reconstruction.  When ACS data is sufficient ($\geq 48$ lines), ESPIRiT achieves 85--95\% of the oracle correction ceiling.  For data-limited regimes ($< 32$ ACS lines) where ESPIRiT degrades, the fallback is a grid search over per-coil amplitude and phase scaling (beam-search); this fallback is also the Sc.~IV method for all modalities without an established specialist calibration algorithm (CASSI, CACTI, compressive holography, ultrasound).  The protocol evaluates three conditions---Sc.~I (true maps), Sc.~II (5\% mismatched maps), and Sc.~IV (ESPIRiT-calibrated maps)---on the same CG-SENSE solver to isolate calibration quality from solver design.

\paragraph{Ultrasound configuration.}
Pulse-echo B-mode imaging is simulated with a 64-element linear array,
5\,MHz center frequency, element pitch $0.3$\,mm, sampling rate 100\,MHz
with 1024 time samples, and speed of sound (SoS) $c = 1540$\,m/s.  The
tissue phantom is a $128 \times 128$ reflectivity map with
Rayleigh-distributed speckle, point scatterers, and anechoic cyst regions.
Frequency-dependent attenuation follows
$\alpha(f) = 0.5\,\text{dB\,cm}^{-1}\text{MHz}^{-1}$.  Mismatch is
parameterized as SoS error $\Delta c \in \{50, 100, 200, 400\}$\,m/s,
perturbing time-of-flight calculations.  The default solver is delay-and-sum
(DAS) beamforming (operator adjoint).  Calibration uses grid search over
21 SoS candidates in $[1440, 1640]$\,m/s, selecting the value that
maximizes reconstruction PSNR.

\paragraph{Controlled hardware experiment protocol.}
The software-perturbation protocol above applies calibrated mask shifts to
existing real measurements.  A full hardware-in-the-loop validation
requires physically displacing the coded aperture mask and re-acquiring
data.  The protocol proceeds as follows: (i)~acquire a baseline dataset
with the mask at its factory-calibrated position; (ii)~physically translate
the mask by a known displacement ($\Delta x \in \{0.25, 0.5, 1.0\}$\,px
equivalent, verified by micrometer stage) and re-acquire under identical
illumination; (iii)~reconstruct both datasets with the factory mask
specification and compute the PSNR degradation and measurement residual;
(iv)~apply \pwm{} autonomous calibration and measure recovery.  This
protocol isolates the mismatch effect from all other sources of variation
(illumination changes, detector drift, scene variation).  Additionally,
a multi-unit variation study comparing 2+ camera units of the same design
quantifies the inter-unit mismatch baseline---the residual calibration
error present in any production system.  Clinical CT phantom validation
(ACR Gammex 464) and clinical MRI validation (fastMRI~\cite{zbontar2018fastmri})
configurations are detailed in Supplementary Note~S25.

\paragraph{Real experimental data.}
Three of the five multi-phantom modalities are validated on real
experimental data from established benchmarks: ultrasound uses PICMUS
experimental phantom recordings and in-vivo carotid
scans~\cite{liebgott2016picmus} plus one DeepUS CIRS-040GSE
acquisition~\cite{khan2020deepus}; cryo-EM uses five EMDB 3D density maps
(EMD-5778, EMD-2984, EMD-6287, EMD-11103, EMD-21375) projected at random
orientations; CT uses three LoDoPaB-CT patient
slices~\cite{leuschner2021lodopab}, one FIPS walnut \textmu CT
scan~\cite{hamalainen2015fips}, and one HTC\,2022
sample~\cite{meaney2022htc}.  Compressive holography and fluorescence
microscopy retain synthetic multi-phantom benchmarks as no public datasets
with calibrated ground truth are available for these modalities.

%% =========================================================================
\subsection*{Statistical Analysis}
\label{sec:methods:statistics}

\paragraph{Per-scene reporting.}
All metrics are reported per scene, not merely as dataset averages.  This
enables identification of scene-dependent failure modes (\eg{}, spectrally
flat scenes that are inherently harder for CASSI, or textureless regions
that challenge SPC).

\paragraph{Summary statistics.}
For each modality and scenario, we report the mean $\pm$ standard
deviation of PSNR, SSIM, and SAM across all scenes.  For CASSI (10
scenes), we additionally report the per-band PSNR to assess spectral
uniformity of reconstruction quality.

\paragraph{Recovery ratio confidence intervals.}
The recovery ratio $\recoveryratio$ (\cref{eq:recovery_ratio}) is a ratio
of differences and therefore sensitive to noise in the constituent PSNR
values.  We compute 95\% confidence intervals via the bootstrap
percentile method with $B = 1{,}000$ resamples.  At each bootstrap
iteration, we resample the scene set with replacement, recompute the
mean PSNR for each scenario, and derive $\recoveryratio$.  The 2.5th
and 97.5th percentiles of the bootstrap distribution define the 95\% CI.

\paragraph{Parameter recovery accuracy.}
For mismatch correction experiments, we report the root-mean-square error
(RMSE) between the estimated and true mismatch parameters:
\begin{equation}
  \text{RMSE}_k = \sqrt{\frac{1}{N_{\text{scene}}}
    \sum_{i=1}^{N_{\text{scene}}}
    (\hat\theta_{k,i} - \theta_{k,\text{true}})^2}
\end{equation}
where $k$ indexes the mismatch parameter, $i$ indexes the scene, and
$N_{\text{scene}}$ is the number of test scenes.  Uncertainty in the RMSE
is estimated via bootstrap ($B = 1{,}000$).

\paragraph{Ablation significance.}
Ablation studies (removal of the photon-budget, recoverability, or
mismatch scoring module, or RunBundle discipline) are evaluated by comparing the
full-pipeline PSNR against each ablated variant.  We report the PSNR
difference $\Delta\psnr$ per modality and verify that each component
contributes $\geq 0.5$\,dB across all validated modalities, establishing
practical significance.


%% =========================================================================
\subsection*{Code and Data Availability}
\label{sec:methods:availability}

\paragraph{Source code.}
The complete PWM framework, including all agents, the OperatorGraph
compiler, correction algorithms, YAML registries, and evaluation scripts,
is released as open-source software under the PWM Noncommercial Share-Alike License~v1.0 at
\url{https://github.com/integritynoble/Physics_World_Model}.  The
codebase is organized into two Python packages:
\texttt{pwm\_core} (core framework, agents, graph compiler, calibration
algorithms) and \texttt{pwm\_AI\_Scientist} (automated experiment
generation and analysis).

\paragraph{Reconstruction data.}
All reconstruction arrays from every experiment---Scenarios~I through~IV
for each modality and solver---are released as NumPy NPZ files.
Files are stored using Git LFS and require \texttt{allow\_pickle=True}
for loading.  Data ranges are standardized: CASSI and CACTI
reconstructions are normalized to $[0, 1]$; SPC reconstructions are in
$[0, 255]$.

\paragraph{Experiment manifests.}
Every experiment is recorded in a RunBundle v0.3.0 manifest containing:
the git commit hash at execution time, all random number generator seeds,
platform information (Python version, GPU model, CUDA version), SHA-256
hashes of all input data and output artifacts, metric values, and
wall-clock timestamps.  These manifests enable exact reproduction of every
reported result.

\paragraph{Registry data.}
All 9 YAML registries (7,034 lines total) that drive the agent system---including
modality definitions, graph templates, photon models, mismatch databases,
compression tables, solver configurations, primitive specifications,
dataset paths, and acceptance thresholds---are publicly available in the
repository under \texttt{packages/pwm\_core/contrib/}.
The \texttt{ExperimentSpec} JSON schemas used for pipeline input
validation are included alongside worked examples in
\texttt{examples/}.
 from the main manuscript.

\section*{Online Methods}

%% =========================================================================
\subsection*{OperatorGraph Specification}
\label{sec:methods:operatorgraph}

\paragraph{Formal definition.}
The \opg{} intermediate representation encodes the forward physics of any
computational imaging modality as a directed acyclic graph (DAG)
$\mathcal{G} = (\mathcal{V}, \mathcal{E})$.
Each node $v_i \in \mathcal{V}$ wraps a \emph{primitive operator} and
implements two entry points: $\texttt{forward}(x) \to y$ and
$\texttt{adjoint}(y) \to x$, the latter defined only when the primitive
is linear.  Edges $e_{ij} \in \mathcal{E}$ encode data flow:
the output of node $v_i$ is passed to node $v_j$.
Each node additionally exposes a set of learnable parameters
$\boldsymbol{\theta}_i$ that may be perturbed during mismatch simulation
or optimized during calibration, as well as read-only metadata flags
(\texttt{is\_linear}, \texttt{is\_stochastic}, \texttt{is\_differentiable}).
The graph is stored as a declarative YAML specification
(\texttt{OperatorGraphSpec}) and compiled to an executable
\texttt{GraphOperator} object by the \texttt{GraphCompiler}.

\paragraph{Node types.}
Primitive operators fall into two categories:

\begin{itemize}
  \item \textbf{Linear operators.}  Convolution (\texttt{conv2d}), mask
    modulation (\texttt{mask\_modulate}), sub-pixel shift (\texttt{subpixel\_shift\_2d}),
    Radon transform (\texttt{radon\_fanbeam}),
    Fourier encoding (\texttt{fourier\_encode}), spectral dispersion
    (\texttt{spectral\_disperse}), Fresnel propagation
    (\texttt{fresnel\_propagate}), random projection
    (\texttt{random\_project}), and structured illumination
    (\texttt{sim\_modulate}).
    Each implements both \texttt{forward()} and \texttt{adjoint()}.
  \item \textbf{Nonlinear operators.}  Squared magnitude
    (\texttt{magnitude\_sq}), Poisson--Gaussian noise
    (\texttt{poisson\_gaussian}), saturation clipping
    (\texttt{saturation\_clip}), phase retrieval nonlinearity
    (\texttt{phase\_abs}), and detector quantization
    (\texttt{quantize}).
    These set \texttt{is\_linear = False} and raise
    \texttt{NotImplementedError} on \texttt{adjoint()}, except where a
    well-defined pseudo-adjoint exists (\eg{}, the identity adjoint for
    magnitude-squared in Gerchberg--Saxton-type algorithms).
\end{itemize}

\paragraph{Adjoint consistency and Lipschitz bounds as mathematical safety rails.}
Two structural constraints distinguish the \opg{} IR from standard deep-learning
forward models and prevent either nonlinear primitive family from becoming a
universal approximator.

First, \emph{adjoint consistency}: correctness of every linear primitive
is verified by a randomized dot-product test.  For a primitive $A$ with
forward map $A:\mathbb{R}^n \to \mathbb{R}^m$, we draw $x \sim \mathcal{N}(0,I_n)$
and $y \sim \mathcal{N}(0,I_m)$ and compute
\begin{equation}
  \delta = \frac{|\langle A^* y,\, x \rangle - \langle y,\, Ax \rangle|}
               {\max(|\langle A^* y,\, x \rangle|,\; \epsilon)}
  \label{eq:adjoint_check}
\end{equation}
where $\epsilon = 10^{-12}$ guards against division by zero.  The test
is repeated $n_{\text{trials}} = 5$ times with independent random draws;
the primitive passes if $\delta_{\max} < 10^{-6}$.  At the graph level,
a compiled \texttt{GraphOperator} composed entirely of linear nodes
executes the same test over the composed forward--adjoint chain.  A
\texttt{GraphAdjointCheckReport} records $n_{\text{trials}}$,
$\delta_{\max}$, and $\bar\delta$ for audit.  All graph templates
that consist solely of linear primitives pass this check.
Adjoint consistency is not merely a numerical correctness guarantee: it
is a \emph{physical faithfulness certificate}.  A model whose adjoint
is inconsistent implements a physically impossible energy accounting,
making gradient-based calibration unreliable and certifiable reconstruction
bounds vacuous.  Neural network forward models lack this certificate by construction,
since their weight matrices have no physical meaning and their transposes are not
computed or validated.

Second, \emph{Lipschitz bounds}: the two nonlinear primitive families,
Detect~$D$ and Transform~$\Lambda$, are each restricted to five canonical
function families with at most two scalar parameters.  This restriction
enforces bounded Lipschitz constants ($\mathrm{Lip}(\Lambda_k) \leq L$ for
the explicit $L$ derived in~\cite{yang2026fpt}) and prevents either primitive
from becoming a universal approximator in the sense of Cybenko~\cite{cybenko1989approximation}.
Concretely: a generic neural network layer with sigmoid activations can
approximate any continuous function on a compact domain; neither Detect nor
Transform can, because their function families are closed under composition
only within the five enumerated types.  This non-universality is the key
property that makes the 11-primitive library \emph{falsifiable}: if a
modality's forward physics cannot be represented within the restricted
families, the extension protocol (Main Text, ``Extension protocol'') fires
and a new primitive must be justified physically.  The combination of
adjoint consistency and bounded Lipschitz nonlinearity provides the
mathematical safety rails that make \opg{}-based reconstruction both
interpretable and certifiable---properties that standard physics-informed
neural networks~\cite{raissi2019pinn} achieve only approximately and
without modality-agnostic guarantees.

\paragraph{Graph compilation.}
The compiler executes a four-stage pipeline:
\begin{enumerate}
  \item \textbf{Validate.}  Confirm acyclicity via topological sort
    (Kahn's algorithm), verify that every \texttt{primitive\_id}
    exists in the global \texttt{PRIMITIVE\_REGISTRY}, reject
    duplicate \texttt{node\_id} values, and optionally verify shape
    compatibility along edges when a \texttt{canonical\_chain}
    metadata flag is set.
  \item \textbf{Bind.}  Instantiate each primitive with its parameter
    dictionary $\boldsymbol{\theta}_i$.
  \item \textbf{Plan forward.}  The topological sort yields a sequential
    execution plan $(v_{\pi(1)}, \ldots, v_{\pi(|\mathcal{V}|)})$.
  \item \textbf{Plan adjoint.}  For graphs where
    \texttt{all\_linear = True}, the adjoint plan reverses the
    topological order and applies each node's individual adjoint in
    sequence, implementing the chain rule $A^* = A_1^*
    \circ \cdots \circ A_{|\mathcal{V}|}^*$ for a composition $A = A_{|\mathcal{V}|}
    \circ \cdots \circ A_1$.  For graphs containing nonlinear nodes,
    the adjoint plan is not generated, and any call to
    \texttt{adjoint()} raises \texttt{NotImplementedError} at runtime.
\end{enumerate}
The compiled \texttt{GraphOperator} is serializable to JSON and hashable
via SHA-256 for provenance tracking in RunBundle manifests.

\paragraph{Template library.}
The \texttt{graph\_templates.yaml} registry contains 170 templates
spanning all five carrier families. Of these, 12 modalities carry full
end-to-end Scenario~I--IV correction validation across all five carrier
families (Supplementary Table~S3). Representative modalities by carrier family include:

\begin{itemize}
  \item \textbf{Optical and X-ray photons:}  CASSI, SPC, CACTI, structured
    illumination microscopy (SIM), confocal, light-sheet, holography,
    ptychography, Fourier ptychographic microscopy (FPM), optical coherence
    tomography (OCT), lensless imaging, light field, integral imaging,
    neural radiance fields (NeRF), Gaussian splatting, fluorescence
    lifetime imaging (FLIM), diffuse optical tomography (DOT),
    phase retrieval, X-ray computed tomography (CT), and cone-beam CT
    (CBCT).
  \item \textbf{Electrons:}  Electron diffraction, electron backscatter
    diffraction (EBSD), electron energy loss spectroscopy (EELS), and
    electron holography.
  \item \textbf{Nuclear spins (MRI):}  Functional MRI (fMRI), diffusion-weighted
    MRI (DW-MRI), and magnetic resonance spectroscopy (MRS).
  \item \textbf{Acoustic waves:}  Ultrasound B-mode, Doppler ultrasound,
    shear-wave elastography, sonar, and photoacoustic tomography
    (combines optical excitation with acoustic detection).
  \item \textbf{Particle-beam probes:}  Neutron tomography, proton CT,
    and muon tomography.  Each decomposes into existing primitives with
    no extension required: neutron CT follows the X-ray CT DAG
    ($P \!\to\! \Lambda \!\to\! D$, Beer--Lambert attenuation with nuclear
    cross-section $\mu_n$); proton CT inserts a Bethe--Bloch
    stopping-power stage as Transform~$\Lambda$ (\#11); muon tomography
    uses Scatter~$R$ (\#10) for multiple Coulomb scattering and
    POCA vertex reconstruction (Supplementary Note~S24).
\end{itemize}

\paragraph{Fidelity levels.}
Each template is parameterized by a fidelity level that controls the
degree of physical realism in the simulated forward model:

\begin{description}
  \item[Level~1 (Linear, shift-invariant):]  The forward model is a
    linear, spatially uniform operator---the simplest approximation,
    suitable for initial diagnostics and rapid prototyping.
  \item[Level~2 (Linear, shift-variant):]  Spatially varying operator
    parameters (e.g.\ non-uniform illumination, position-dependent PSF,
    multi-coil sensitivity maps in MRI).  Adds a modality-appropriate
    noise model (Poisson shot noise plus Gaussian read noise for
    photon-counting modalities, Rician noise for MRI, Poisson for CT).
  \item[Level~3 (Nonlinear, ray/wave-based):]  Includes nonlinear
    effects such as beam hardening, phase wrapping, and scattering,
    captured by the Transform~$\Lambda$ primitive.
    Perturbation families and ranges are specified in
    \texttt{mismatch\_db.yaml}.
  \item[Level~4 (Full-wave / Monte Carlo):]  Complete physical
    simulation including wave-optical propagation, spatially varying
    aberrations, detector nonlinearities, and environmental drift.
    Currently implemented for holography and ptychography.
\end{description}

\noindent All 11 primitives in the library $\Blib$ apply uniformly across every fidelity level; the levels organize simulation complexity, not theorem scope.


%% =========================================================================
\subsection*{Triad Decomposition Formalization}
\label{sec:methods:triad}

The \triad{} asserts that the quality of any computational imaging
reconstruction is bounded by three fundamental gates.  Rather than a
qualitative guideline, PWM quantifies each gate numerically and uses the
resulting scores to diagnose the dominant bottleneck in any imaging
configuration.

\paragraph{Gate~1 (Recoverability).}
Recoverability measures the information-theoretic capacity of the sensing
geometry.  We quantify it via the \emph{effective compression ratio}
$r = m / n$, where $m$ is the number of independent measurements and $n$
the dimension of the signal.  The \texttt{compression\_db.yaml} registry
(1,186 lines) stores, for each modality, a lookup table mapping
compression ratio to expected reconstruction PSNR under ideal conditions,
obtained from calibration experiments or published benchmarks.  Each entry
carries a \texttt{provenance} field citing the source (paper DOI, internal
experiment ID, or theoretical formula).  Additional recoverability
indicators include the effective rank of the measurement matrix (estimated
via randomized SVD for large operators), the dimension of the null space,
and the restricted isometry property (RIP) constant where analytically
tractable (\eg{}, for Gaussian random projections in SPC).

\paragraph{Gate~2 (Carrier Budget).}
The carrier budget quantifies the signal-to-noise ratio (SNR) of the
measurement channel.  The photon-budget module consumes the
\texttt{photon\_db.yaml} registry which stores, per modality,
a deterministic photon model parameterized by source power, quantum
efficiency, exposure time, and detector characteristics.  It
classifies the noise regime into one of three categories:
\emph{shot-limited} (Poisson-dominated, SNR $\propto \sqrt{N_{\text{photon}}}$),
\emph{read-limited} (Gaussian read noise dominates, SNR $\propto
N_{\text{photon}} / \sigma_{\text{read}}$), and \emph{dark-current-limited}
(long exposures where dark current accumulation dominates).  The output is
a \texttt{PhotonReport} containing the estimated SNR in decibels, the noise
regime classification, per-element photon count, and a feasibility verdict
(\texttt{sufficient}, \texttt{marginal}, or \texttt{insufficient}).

\paragraph{Gate~3 (Operator Mismatch).}
Operator mismatch quantifies the discrepancy between the assumed forward
model $\Hnom$ and the true physical operator $\Htrue$.  The
mismatch scoring module consults \texttt{mismatch\_db.yaml}
which catalogs, for each modality, the set of mismatch parameters
(spatial shifts, rotational offsets, dispersion errors, PSF deviations,
coil sensitivity errors, center-of-rotation offsets, \etc{}), their
typical ranges, and available correction methods.  The mismatch severity
score $s \in [0, 1]$ is computed as the normalized $\ell_2$ distance
$\| \boldsymbol{\theta}_{\text{true}} - \boldsymbol{\theta}_{\text{nom}} \|
/ \| \boldsymbol{\theta}_{\text{range}} \|$, where
$\boldsymbol{\theta}_{\text{range}}$ is the per-parameter dynamic range
from the registry.  Sensitivity analysis $\partial \psnr / \partial
\theta_k$ is estimated via finite differences on the forward model.  The
output is a \texttt{MismatchReport} containing the severity score, the
dominant mismatch parameter, the recommended correction method, and the
expected PSNR gain from correction.

\paragraph{Gate binding determination.}
Given reconstruction results under the four-scenario protocol
(the Evaluation Protocol section below), PWM identifies the dominant gate by
comparing three cost terms:
\begin{align}
  C_{\text{mismatch}} &= \psnr_{\text{I}} - \psnr_{\text{II}} \label{eq:cost_mismatch} \\
  C_{\text{noise}}    &= \psnr_{\text{ideal}} - \psnr_{\text{noisy}} \label{eq:cost_noise} \\
  C_{\text{recover}}  &= \psnr_{\text{limit}} - \psnr_{\text{I}} \label{eq:cost_recover}
\end{align}
where $\psnr_{\text{I}}$ is the reconstruction PSNR under Scenario~I
(ideal operator), $\psnr_{\text{II}}$ under Scenario~II (mismatched
operator), $\psnr_{\text{noisy}}$ under the corresponding noisy condition,
and $\psnr_{\text{limit}}$ is the theoretical upper bound from the
compression table.  The dominant gate is $\arg\max_g C_g$.

\paragraph{TriadReport.}
For every diagnosis, \pwm{} produces a \triadreport{}: a Pydantic-validated
structured artifact comprising
\texttt{dominant\_gate} (enum: \texttt{recoverability},
\texttt{carrier\_budget}, \texttt{operator\_mismatch}),
\texttt{evidence\_scores} (three floats, one per gate),
\texttt{confidence\_interval} (float, 95\% CI width from bootstrap),
\texttt{recommended\_action} (string, \eg{} ``increase compression ratio''
or ``apply mismatch correction''), and
\texttt{parameter\_sensitivities} (dictionary mapping each mismatch
parameter name to its $\partial \psnr / \partial \theta_k$ value).
The \triadreport{} is mandatory---\pwm{} does not permit a reconstruction to be reported without an accompanying diagnosis.

\paragraph{Recovery ratio.}
We define the \emph{recovery ratio}
\begin{equation}
  \recoveryratio = \frac{\psnr_{\text{IV}} - \psnr_{\text{II}}}
                        {\psnr_{\text{I}} - \psnr_{\text{II}}}
  \label{eq:recovery_ratio}
\end{equation}
which lies in $[0, 1]$ under standard convexity conditions (see Supplementary Note~1 for formal analysis; values $\recoveryratio > 1$ are possible when the corrected operator provides beneficial regularization).
$\recoveryratio = 0$ indicates that calibration yields no benefit (mismatch
is not the bottleneck), while $\recoveryratio = 1$ indicates that
calibration fully closes the mismatch gap.


%% =========================================================================
\subsection*{Agent System Architecture}
\label{sec:methods:agents}

The three gates are implemented as a deterministic scoring pipeline:
an orchestrator maps the user request to a modality from the registry,
then dispatches three parallel scoring modules---one per gate---whose
reports are fused by a bottleneck classifier that identifies the dominant
gate and by a negotiator that enforces cross-gate veto rules
(e.g., rejecting reconstruction when photon starvation coincides with
aggressive compression).  All quantitative decisions are deterministic;
an LLM is used only as an optional fallback for natural-language modality
mapping and is never on the scoring path.  The full agent architecture
and ablation studies are described in a companion paper.

\paragraph{Contract system.}
Inter-agent communication uses 25 Pydantic v2 contract models.  All
contracts inherit from \texttt{StrictBaseModel}, which enforces
\texttt{extra="forbid"} (no unexpected fields),
\texttt{validate\_assignment=True} (mutations re-validated), and a model
validator that rejects NaN and Inf in any float field.  Bounded scores use
\texttt{Field(ge=0.0, le=1.0)}.  Enums are string enums for human-readable
JSON serialization.  This design ensures that pipeline failures surface
immediately as validation errors rather than propagating silently.

\paragraph{YAML registries.}
The system is driven by 9 YAML registries totalling 7,034 lines:
\texttt{modalities.yaml} (modality definitions),
\texttt{graph\_templates.yaml} (OperatorGraph skeletons),
\texttt{photon\_db.yaml} (photon models),
\texttt{mismatch\_db.yaml} (mismatch parameters and correction methods),
\texttt{compression\_db.yaml} (recoverability tables with provenance),
\texttt{solver\_registry.yaml} (solver configurations),
\texttt{primitives.yaml} (primitive operator metadata),
\texttt{dataset\_registry.yaml} (dataset locations and formats), and
\texttt{acceptance\_thresholds.yaml} (pass/fail thresholds per metric).


%% =========================================================================
\subsection*{Correction Algorithms}
\label{sec:methods:correction}

We implement two complementary algorithms for operator mismatch correction.
Crucially, both algorithms operate on the forward operator parameters
$\boldsymbol{\theta}$ rather than the reconstruction solver weights,
making them \emph{solver-agnostic}: the corrected operator
$H(\hat{\boldsymbol{\theta}})$ benefits any downstream solver (GAP-TV,
MST-L, HDNet~\cite{hu2022hdnet}, CST, \etc{}) without retraining.

\paragraph{Algorithm~1: Hierarchical Beam Search.}
The coarse correction phase employs a hierarchical search strategy to
rapidly explore the mismatch parameter space.  For CASSI, the five-parameter
mismatch model comprises mask affine parameters (spatial shifts $dx$, $dy$
and rotation $\theta$) and dispersion parameters (slope $a_1$ and axis
angle $\alpha$); an optional sixth parameter, PSF width $\sigma_{\text{psf}}$, is available but not used in the primary experiments.
The algorithm proceeds as follows:

\begin{enumerate}
  \item \textbf{1D sweeps.}  Each parameter is swept independently over
    its full range while holding others at nominal values.  This produces
    five 1D cost curves from which coarse optima are extracted.
  \item \textbf{3D beam search.}  The mask affine subspace
    $(dx, dy, \theta)$ is searched over a $5 \times 5 \times 5$ grid
    centered on the 1D optima.  The top-$k$ ($k = 5$) candidates by
    reconstruction PSNR are retained.
  \item \textbf{2D beam search.}  For each retained mask candidate,
    the dispersion subspace $(a_1, \alpha)$ is searched over a
    $5 \times 7$ grid.  The joint top-$k$ candidates are retained.
  \item \textbf{Coordinate descent refinement.}  Three rounds of
    univariate refinement on each parameter, shrinking the search
    interval by factor~2 at each round, produce the final estimate
    $\hat{\boldsymbol{\theta}}_{\text{Alg1}}$.
\end{enumerate}

Total runtime is approximately 300 seconds per scene on a single GPU.
Accuracy is $\pm 0.1$--$0.2$ pixels for spatial parameters and
$\pm 0.05^\circ$ for angular parameters.

\paragraph{Algorithm~2: Joint Gradient Refinement.}
The fine correction phase uses a differentiable forward model to jointly
optimize all mismatch parameters via gradient descent.  The key components
are:

\begin{enumerate}
  \item \textbf{Differentiable mask warp.}  The binary mask is warped by a
    continuous affine transformation using bilinear interpolation,
    implemented as a custom PyTorch module
    (\texttt{DifferentiableMaskWarpFixed}).  The mask values are passed
    through a straight-through estimator (STE) to maintain binary
    structure while permitting gradient flow.
  \item \textbf{Differentiable forward model.}  The CASSI forward model
    $y = \texttt{CASSI}(x;\, \boldsymbol{\theta})$ is implemented as a
    differentiable PyTorch module (\texttt{DifferentiableCassiForwardSTE})
    that accepts mismatch parameters as differentiable inputs.
  \item \textbf{GPU grid initialization.}  A full-range 3D grid search
    over $(dx, dy, \theta)$ with $9 \times 9 \times 7 = 567$ points
    provides diverse starting candidates.  The top 9 candidates seed
    multi-start gradient refinement.
  \item \textbf{Staged gradient refinement.}  Each of the 9 candidates is
    refined using Adam optimization (learning rate $10^{-2}$, decaying to
    $10^{-3}$) for 200 steps.  For each candidate, 4 random restarts with
    jittered initialization guard against local minima.  The loss function
    is the negative PSNR computed via an unrolled $K$-iteration
    differentiable GAP-TV solver (\texttt{DifferentiableGAPTV}, $K=10$
    unrolled iterations).
\end{enumerate}

Total runtime for Algorithm~2 is approximately 3,200 seconds
(200~steps $\times$ 4~restarts $\times$ 9~candidates with early stopping).
Accuracy improves to $\pm 0.05$--$0.1$ pixels, a 3--5$\times$ improvement
over Algorithm~1.  The two algorithms are used sequentially in practice:
Algorithm~1 provides a warm start, and Algorithm~2 refines to sub-pixel
precision.


%% =========================================================================
\subsection*{Evaluation Protocol}
\label{sec:methods:eval}

\paragraph{Four-Scenario Protocol.}
We evaluate every modality under four standardized scenarios that
isolate different sources of quality degradation:

\begin{description}
  \item[Scenario~I (Ideal):]
    $\yobs = \Htrue\, \xgt$; reconstruct with $\Htrue$.
    In this scenario the system is perfectly calibrated ($\Htrue = \Hnom$),
    so the operator used for reconstruction matches the one that generated
    the data.  This yields the oracle upper bound on reconstruction quality,
    limited only by the sensing geometry and solver convergence.
  \item[Scenario~II (Mismatch):]
    $\yobs = \Htrue\, \xgt$; reconstruct with $\Hnom$
    ($\Hnom \neq \Htrue$).
    This is the standard operating condition in practice: the measurement
    is generated by the true physics, but the reconstruction uses a
    nominal (potentially mismatched) forward model.
  \item[Scenario~III (Oracle):]
    Same measurements as Scenario~II
    ($\yobs = \Htrue\, \xgt$ with $\Htrue \neq \Hnom$);
    reconstruct with $\Htrue$ instead of $\Hnom$.
    Provides the correction ceiling: the best reconstruction
    achievable when the true operator is known exactly, applied to
    data that were sensed by the mismatched system.  The gap between
    Scenario~III and Scenario~I reveals the irreducible loss from
    the degraded sensing configuration itself (e.g., a shifted mask
    pattern is suboptimal even when perfectly known).
  \item[Scenario~IV (Corrected):]
    $\yobs = \Htrue\, \xgt$; reconstruct with
    $\hat{H} = H(\hat{\boldsymbol{\theta}})$ where
    $\hat{\boldsymbol{\theta}}$ is estimated by Algorithms~1 and~2.
    This quantifies the benefit of mismatch calibration.
\end{description}

\paragraph{Metrics.}
Reconstruction quality is assessed using three complementary metrics:
\begin{itemize}
  \item \textbf{PSNR} (peak signal-to-noise ratio, in dB): the primary
    metric, computed per scene and averaged.  For signals normalized to
    $[0, 1]$, $\psnr = 10 \log_{10}(1 / \text{MSE})$.  For SPC data
    normalized to $[0, 255]$, the peak value is 255.
  \item \textbf{SSIM} (structural similarity index): captures perceptual
    quality including luminance, contrast, and structural components,
    computed with a Gaussian window of width 11 and standard deviation
    1.5.
  \item \textbf{SAM} (spectral angle mapper): for hyperspectral
    modalities (CASSI), measures the angle between predicted and true
    spectral vectors at each spatial location, reported in degrees.
    Lower is better.
\end{itemize}

\paragraph{Datasets.}
\begin{itemize}
  \item \textbf{CASSI:} 10 scenes from the KAIST dataset~\cite{choi2017kaist},
    each a $256 \times 256 \times 28$ spectral cube (28 spectral bands
    from 450\,nm to 650\,nm).  Data range $[0, 1]$.
  \item \textbf{CACTI:} 6 benchmark videos, each $256 \times 256 \times 8$
    (8 temporal frames encoded per snapshot).  Data range $[0, 1]$.
  \item \textbf{SPC:} 11 natural images from the Set11 benchmark, each
    $256 \times 256$ grayscale.  Data range $[0, 255]$.
\end{itemize}

\noindent All per-scene metrics are reported individually as well as
averaged, and all reconstruction arrays are saved as NumPy NPZ files.


%% =========================================================================
\subsection*{Experimental Details}
\label{sec:methods:experimental}

\paragraph{Hardware.}
All experiments are conducted on a single NVIDIA GPU.  Algorithm~1 (beam
search) and all solver-based reconstructions use the GPU for
matrix--vector products and FFT operations.  Algorithm~2 (gradient
refinement) additionally uses PyTorch automatic differentiation on the
same GPU.

\paragraph{CASSI configuration.}
The coded aperture snapshot spectral imaging (CASSI) system uses a TSA-Net
binary mask of dimensions $256 \times 256$, with 28 spectral bands
dispersed along the spatial dimension.  The five-parameter mismatch model
$\mismatchparam = (dx, dy, \theta, a_1, \alpha)$
describes: mask spatial shift in $x$ ($dx$, pixels), mask spatial shift in
$y$ ($dy$, pixels), mask rotation angle ($\theta$, degrees), dispersion
slope ($a_1$, pixels per band), and dispersion axis angle ($\alpha$,
degrees).  An optional sixth parameter, PSF blur width ($\sigma_{\text{psf}}$, pixels), is available but not used in the primary experiments.
These mismatch parameter values were determined through systematic characterization of realistic CASSI assembly errors (Supplementary Note~8). The true mismatch parameters are
$\mismatchparam_{\text{true}} = (dx = 0.5\;\text{px},\; dy = 0.3\;\text{px},\;
\theta = 0.1^\circ,\; a_1 = 2.02,\;
\alpha = 0.15^\circ)$.
Solvers evaluated include TwIST~\cite{bioucasdias2007twist}, GAP-TV~\cite{yuan2016gaptv}, DGSMP~\cite{huang2021dgsmp}, MST-L~\cite{cai2022mst}, and CST-L~\cite{cai2022cst}, all of which receive the same operator and differ only in their reconstruction algorithm. The supplementary per-scene analysis additionally includes DeSCI~\cite{liu2019rank} and HDNet~\cite{hu2022hdnet}.
\textit{Real-data validation.}
The TSA real hyperspectral dataset~\cite{wagadarikar2008cassi} consists of
5 scenes at $660 \times 660$ spatial resolution with 28 spectral bands
and mask-shift step~2.  Four solvers are evaluated: GAP-TV (200 iterations),
HDNet (pre-trained checkpoint, full spatial resolution), MST-S and MST-L
(pre-trained checkpoints, centre-cropped to $256 \times 256$ due to
hardcoded spatial assumptions in the model architecture).  The coded
aperture mask is perturbed by $dx = 0.5$\,px, $dy = 0.3$\,px to simulate
assembly-induced mismatch.  No ground truth is available; quality is
assessed via the normalised measurement residual
$r = \|\mathbf{y} - H\hat{\mathbf{x}}\|^2 / \|\mathbf{y}\|^2$.

\paragraph{CACTI configuration.}
The coded aperture compressive temporal imaging system uses binary
temporal masks of dimensions $256 \times 256$, encoding 8 video frames
into a single snapshot measurement.  Mismatch is parameterized as a
temporal mask timing offset (sub-frame shift).  The default solver is
EfficientSCI~\cite{wang2023efficientsci}.
\textit{Real-data validation.}
The EfficientSCI real temporal dataset~\cite{wang2023efficientsci} consists
of 4 dynamic scenes (duomino, hand, pendulumBall, waterBalloon) at
$512 \times 512$ with compression ratio~10.  The real mask is stored
separately from the measurement data.  Two solvers are evaluated:
GAP-TV (50 iterations) and PnP-FFDNet (50 iterations with FFDNet denoiser).
Mismatch is induced by shifting the mask by $dx = 0.5$\,px, $dy = 0.3$\,px.
Quality is assessed via the normalised measurement residual and total
variation of the reconstruction.

\paragraph{SPC configuration.}
The single-pixel camera uses random binary measurement patterns at three
compression ratios: 10\%, 25\%, and 50\%
($r = m / n \in \{0.10, 0.25, 0.50\}$).  Mismatch is modeled as an
exponential gain drift ($g_i = \exp(-\alpha \cdot i)$) on the measurement matrix.  The default solver
is FISTA-TV with total-variation regularization.

\paragraph{Lensless imaging configuration.}
Lensless cameras replace conventional optics with a diffuser placed directly
on the sensor, making the forward model a 2D convolution with the measured
point-spread function (PSF).  The object is a $64 \times 64$ image
convolved with a Gaussian PSF of width $\sigma = 2.0$\,px.  Mismatch is
parameterized as PSF width error
$\Delta\sigma \in \{0.5, 1.0, 2.0, 3.0\}$\,px, simulating an error in
the assumed pinhole-to-sensor distance.  Five structurally diverse phantoms
are used ($N{=}5$).  The default solver is FISTA-TV (80 iterations).
Calibration uses forward-model fitting on a known reference target:
$\sigma_{\mathrm{cal}} = \mathrm{argmin}_\sigma \|y_{\mathrm{cal}} -
h_\sigma * x_{\mathrm{cal}}\|^2$.
\textit{Real-data validation.}
Five DiffuserCam DLMD Mirflickr scenes~\cite{Antipa2018} are used to
corroborate simulation results on physical photographic content.

\paragraph{Compressive holography configuration.}
Multi-depth inline holography encodes a $(4 \times 64 \times 64)$ 3D
object into a single $(64 \times 64)$ hologram via Fresnel propagation at
four depth planes with spacing $200$\,\textmu m, wavelength
$\lambda = 532$\,nm, and pixel pitch $5$\,\textmu m.  Mismatch is
parameterized as propagation distance error
$\Delta z \in \{10, 50, 100, 200\}$\,\textmu m applied uniformly to all
depth planes.  The default solver is FISTA-TV (80 iterations,
$\lambda_{\mathrm{TV}} = 0.005$).  Calibration uses hologram residual
minimisation: grid search over distance error candidates with the search
range clamped to ensure all propagation distances remain positive.

\paragraph{Fluorescence microscopy configuration.}
Widefield fluorescence imaging uses a dual-PSF Stokes-shift model with
Gaussian excitation PSF ($\sigma_{\mathrm{ex}} = 1.5$\,px) and emission
PSF ($\sigma_{\mathrm{em}} = 2.0$\,px), quantum yield $\eta = 0.7$, and
background level $b = 0.02$.  The specimen is a $64 \times 64$ image with
fluorescent puncta, diffuse organelles, and filamentary structures.
Mismatch is parameterized as PSF sigma error
$\Delta\sigma \in \{0.1, 0.3, 0.5, 1.0\}$\,px applied to both excitation
and emission PSFs.  The default solver is Richardson--Lucy (80 iterations).
Calibration uses a $9 \times 9$ grid search over
$(\sigma_{\mathrm{ex}}, \sigma_{\mathrm{em}})$ centered at the nominal
values, selecting the pair that minimizes the measurement residual.

\paragraph{CT configuration.}
Fan-beam geometry with 180 projections over $180^\circ$.  Mismatch is
modeled as a center-of-rotation (CoR) offset, which produces
characteristic arc artifacts in the reconstruction.  The default solver
is filtered back-projection (FBP)~\cite{feldkamp1984practical} with a Ram-Lak filter, supplemented
by iterative SART for comparison.

\paragraph{CT detector offset configuration.}
Parallel-beam CT is simulated with 360 projection angles over $[0, \pi)$
and 128 detector elements.  The object is a $128 \times 128$ Shepp--Logan
phantom.  Detector offset mismatch is simulated by shifting the sinogram
along the detector axis by $\Delta s \in \{2, 5, 10, 20\}$\,pixels.  The
default solver is filtered backprojection (FBP) with Shepp--Logan filter.
Calibration uses grid search over 21 offset candidates in $[-25, 25]$\,px,
selecting the value that maximizes reconstruction PSNR.

\paragraph{Ptychography configuration.}
Electron ptychography uses 4D-STEM data: a focused probe is scanned across
a $128 \times 128$ grid, recording a convergent-beam electron diffraction
(CBED) pattern at each position (300\,kV, SrTiO$_3$ [001] from
Zenodo~5113449).  The forward model is Fresnel near-field propagation of
the exit wave through the specimen.  Mismatch is parameterized as probe
position jitter $\Delta p \in \{0.5, 1.0, 2.0, 4.0\}$\,px, simulating
scan-coil hysteresis or sample drift.  The default solver is integrated
centre-of-mass (iCoM) phase retrieval.  Calibration uses ePIE blind
position correction~\cite{maiden2009ptychography}.

\paragraph{Cryo-EM configuration.}
Cryo-EM imaging is modeled as contrast transfer function (CTF) filtering
of a 2D projected potential ($128 \times 128$, pixel size $0.1$\,nm).  The
CTF is parameterized by defocus $\Delta f = -2000$\,nm, spherical aberration
$C_s = 2.0$\,mm, electron wavelength $\lambda = 0.00251$\,nm (300\,kV),
with a B-factor envelope ($B = 2$\,nm$^2$) and ice-thickness attenuation
(mean free path $350$\,nm, thickness $50$\,nm).  Mismatch is parameterized
as defocus error $\Delta(\Delta f) \in \{100, 200, 500, 1000\}$\,nm.  The
default solver is Wiener filtering with SNR\,$= 50$.  Calibration uses grid
search over 51 defocus values in $[-3000, -500]$\,nm, selecting the value
that maximizes reconstruction PSNR.

\paragraph{MRI configuration.}
Cartesian $k$-space sampling with 4$\times$ acceleration (25\% of
$k$-space lines acquired).  Mismatch is parameterized as a 5\% multiplicative
error in the coil sensitivity maps used for parallel imaging reconstruction.
The default solver is CG-SENSE~\cite{pruessmann1999sense} with $\ell_1$-wavelet regularization ($\lambda = 10^{-3}$, 30 iterations).
\textit{Scenario~IV calibration (ESPIRiT-adopted).}
For MRI, Scenario~IV adopts ESPIRiT~\cite{uecker2014espirit} as the calibration method: coil sensitivity maps are estimated directly from the central $k$-space autocalibration signal (ACS) region via eigenvalue decomposition of the calibration matrix, then used in the CG-SENSE reconstruction.  When ACS data is sufficient ($\geq 48$ lines), ESPIRiT achieves 85--95\% of the oracle correction ceiling.  For data-limited regimes ($< 32$ ACS lines) where ESPIRiT degrades, the fallback is a grid search over per-coil amplitude and phase scaling (beam-search); this fallback is also the Sc.~IV method for all modalities without an established specialist calibration algorithm (CASSI, CACTI, compressive holography, ultrasound).  The protocol evaluates three conditions---Sc.~I (true maps), Sc.~II (5\% mismatched maps), and Sc.~IV (ESPIRiT-calibrated maps)---on the same CG-SENSE solver to isolate calibration quality from solver design.

\paragraph{Ultrasound configuration.}
Pulse-echo B-mode imaging is simulated with a 64-element linear array,
5\,MHz center frequency, element pitch $0.3$\,mm, sampling rate 100\,MHz
with 1024 time samples, and speed of sound (SoS) $c = 1540$\,m/s.  The
tissue phantom is a $128 \times 128$ reflectivity map with
Rayleigh-distributed speckle, point scatterers, and anechoic cyst regions.
Frequency-dependent attenuation follows
$\alpha(f) = 0.5\,\text{dB\,cm}^{-1}\text{MHz}^{-1}$.  Mismatch is
parameterized as SoS error $\Delta c \in \{50, 100, 200, 400\}$\,m/s,
perturbing time-of-flight calculations.  The default solver is delay-and-sum
(DAS) beamforming (operator adjoint).  Calibration uses grid search over
21 SoS candidates in $[1440, 1640]$\,m/s, selecting the value that
maximizes reconstruction PSNR.

\paragraph{Controlled hardware experiment protocol.}
The software-perturbation protocol above applies calibrated mask shifts to
existing real measurements.  A full hardware-in-the-loop validation
requires physically displacing the coded aperture mask and re-acquiring
data.  The protocol proceeds as follows: (i)~acquire a baseline dataset
with the mask at its factory-calibrated position; (ii)~physically translate
the mask by a known displacement ($\Delta x \in \{0.25, 0.5, 1.0\}$\,px
equivalent, verified by micrometer stage) and re-acquire under identical
illumination; (iii)~reconstruct both datasets with the factory mask
specification and compute the PSNR degradation and measurement residual;
(iv)~apply \pwm{} autonomous calibration and measure recovery.  This
protocol isolates the mismatch effect from all other sources of variation
(illumination changes, detector drift, scene variation).  Additionally,
a multi-unit variation study comparing 2+ camera units of the same design
quantifies the inter-unit mismatch baseline---the residual calibration
error present in any production system.  Clinical CT phantom validation
(ACR Gammex 464) and clinical MRI validation (fastMRI~\cite{zbontar2018fastmri})
configurations are detailed in Supplementary Note~S25.

\paragraph{Real experimental data.}
Three of the five multi-phantom modalities are validated on real
experimental data from established benchmarks: ultrasound uses PICMUS
experimental phantom recordings and in-vivo carotid
scans~\cite{liebgott2016picmus} plus one DeepUS CIRS-040GSE
acquisition~\cite{khan2020deepus}; cryo-EM uses five EMDB 3D density maps
(EMD-5778, EMD-2984, EMD-6287, EMD-11103, EMD-21375) projected at random
orientations; CT uses three LoDoPaB-CT patient
slices~\cite{leuschner2021lodopab}, one FIPS walnut \textmu CT
scan~\cite{hamalainen2015fips}, and one HTC\,2022
sample~\cite{meaney2022htc}.  Compressive holography and fluorescence
microscopy retain synthetic multi-phantom benchmarks as no public datasets
with calibrated ground truth are available for these modalities.

%% =========================================================================
\subsection*{Statistical Analysis}
\label{sec:methods:statistics}

\paragraph{Per-scene reporting.}
All metrics are reported per scene, not merely as dataset averages.  This
enables identification of scene-dependent failure modes (\eg{}, spectrally
flat scenes that are inherently harder for CASSI, or textureless regions
that challenge SPC).

\paragraph{Summary statistics.}
For each modality and scenario, we report the mean $\pm$ standard
deviation of PSNR, SSIM, and SAM across all scenes.  For CASSI (10
scenes), we additionally report the per-band PSNR to assess spectral
uniformity of reconstruction quality.

\paragraph{Recovery ratio confidence intervals.}
The recovery ratio $\recoveryratio$ (\cref{eq:recovery_ratio}) is a ratio
of differences and therefore sensitive to noise in the constituent PSNR
values.  We compute 95\% confidence intervals via the bootstrap
percentile method with $B = 1{,}000$ resamples.  At each bootstrap
iteration, we resample the scene set with replacement, recompute the
mean PSNR for each scenario, and derive $\recoveryratio$.  The 2.5th
and 97.5th percentiles of the bootstrap distribution define the 95\% CI.

\paragraph{Parameter recovery accuracy.}
For mismatch correction experiments, we report the root-mean-square error
(RMSE) between the estimated and true mismatch parameters:
\begin{equation}
  \text{RMSE}_k = \sqrt{\frac{1}{N_{\text{scene}}}
    \sum_{i=1}^{N_{\text{scene}}}
    (\hat\theta_{k,i} - \theta_{k,\text{true}})^2}
\end{equation}
where $k$ indexes the mismatch parameter, $i$ indexes the scene, and
$N_{\text{scene}}$ is the number of test scenes.  Uncertainty in the RMSE
is estimated via bootstrap ($B = 1{,}000$).

\paragraph{Ablation significance.}
Ablation studies (removal of the photon-budget, recoverability, or
mismatch scoring module, or RunBundle discipline) are evaluated by comparing the
full-pipeline PSNR against each ablated variant.  We report the PSNR
difference $\Delta\psnr$ per modality and verify that each component
contributes $\geq 0.5$\,dB across all validated modalities, establishing
practical significance.


%% =========================================================================
\subsection*{Code and Data Availability}
\label{sec:methods:availability}

\paragraph{Source code.}
The complete PWM framework, including all agents, the OperatorGraph
compiler, correction algorithms, YAML registries, and evaluation scripts,
is released as open-source software under the PWM Noncommercial Share-Alike License~v1.0 at
\url{https://github.com/integritynoble/Physics_World_Model}.  The
codebase is organized into two Python packages:
\texttt{pwm\_core} (core framework, agents, graph compiler, calibration
algorithms) and \texttt{pwm\_AI\_Scientist} (automated experiment
generation and analysis).

\paragraph{Reconstruction data.}
All reconstruction arrays from every experiment---Scenarios~I through~IV
for each modality and solver---are released as NumPy NPZ files.
Files are stored using Git LFS and require \texttt{allow\_pickle=True}
for loading.  Data ranges are standardized: CASSI and CACTI
reconstructions are normalized to $[0, 1]$; SPC reconstructions are in
$[0, 255]$.

\paragraph{Experiment manifests.}
Every experiment is recorded in a RunBundle v0.3.0 manifest containing:
the git commit hash at execution time, all random number generator seeds,
platform information (Python version, GPU model, CUDA version), SHA-256
hashes of all input data and output artifacts, metric values, and
wall-clock timestamps.  These manifests enable exact reproduction of every
reported result.

\paragraph{Registry data.}
All 9 YAML registries (7,034 lines total) that drive the agent system---including
modality definitions, graph templates, photon models, mismatch databases,
compression tables, solver configurations, primitive specifications,
dataset paths, and acceptance thresholds---are publicly available in the
repository under \texttt{packages/pwm\_core/contrib/}.
The \texttt{ExperimentSpec} JSON schemas used for pipeline input
validation are included alongside worked examples in
\texttt{examples/}.
